\chapter{Haar Measure}

\section{Reminder of General Topology}

\begin{definitionenv}[Open Mapping]\label{def:IsOpenMap}
    Let $X$ and $Y$ be topological spaces and $f: X\rightarrow Y$ be a mapping.
    We say that $f$ is an open mapping if for any open $U \subseteq X$, $f(U)$ is open in $Y$.
\end{definitionenv}

\begin{definitionenv}[Quotient Topology] % Mathlib uses coinduced topology to define this.
    Let $(X,\mathscr{T})$ be a topological space and $\pi : X\rightarrow Y$ be a surjective mapping.
    We call quotient topology of $\mathscr{T}$ on $Y$ (with respect to $\pi$) the family $V\subseteq Y$ such that 
    $\pi^{-1}(V) \in \mathscr{T}$.

    \quad Note that $\pi$ is continuous if we equipped $Y$ with the quotient topology.
\end{definitionenv}

\begin{propositionenv}
    Let $X,Y$ be topological spaces, $\pi : X\rightarrow Y$ be a surjective mapping.
    We equipped $Y$ with the quotient topology and assume that $\pi$ is open.
    \begin{enumerate}
        \item A subset $U$ of $Y$ is open if and only if $\pi^{-1}(U)$ is open in $X$.
        \item A subset $F$ of $Y$ is closed if and only if $\pi^{-1}(F)$ is closed in $X$.
    \end{enumerate}
\end{propositionenv}
\begin{proofenv}
    \quad 
    \begin{enumerate}
        \item $\pi$ is continuous, so if $U$ is open in $Y$, then $\pi^{-1}(U)$ is open in $X$.
        Conversely, if $\pi^{-1}(U)$ is open in $X$, then $\pi\left(\pi^{-1}(U)\right) = U$ is open.
        \item If $F$ is closed in $Y$, then $Y\setminus F$ is open, hence
        $X\setminus \pi^{-1}(F)=\pi^{-1}(Y\setminus F)$
        is open.
        Conversely, if $X\setminus \pi^{-1}(F)$ is open, then $\pi\left(\pi^{-1}\left(Y\setminus F\right)\right) = Y\setminus F$ is open.
    \end{enumerate}
\end{proofenv}
\begin{propositionenv}\label{prop:5.1.4}
    Let $I$ be a set, $\left(X_i\right)_{i\in I}$, $\left(Y_i\right)_{i\in I}$ be families of topological spaces.
    For any $i\in I$, let $f_i: X_i \rightarrow Y_i$ be a mapping.
    Let
    $$ X = \prod_{i\in I} X_i, \ Y = \prod_{i\in I} Y_i$$
    (Considered with initial topology) and let
    $$\begin{array}{rrcl}
        f : & X & \longrightarrow & Y\\
        & \left(x_i\right)_{i\in I} & \longmapsto & \left(f_i(x_i)\right)_{i\in I}.
    \end{array}$$

    \begin{enumerate}
        \item If $f_i$ is continuous for any $i\in I$, then $f$ is continuous.
        \item If $f_i$ is open for any $i\in I$ and is surjective for all but finitely many $i\in I$, then $f$ is open.
    \end{enumerate}
\end{propositionenv}
\begin{proofenv}
    \quad
    \begin{enumerate}
        \item For any $i\in I$, let $\pi_i : X\rightarrow X_i$, $p_i : Y \rightarrow Y_i$ be projections (all continuous).
            \begin{center}
                \begin{tikzcd}
                	X & Y \\
                	{X_i} & {Y_i}
                	\arrow["f", from=1-1, to=1-2]
                	\arrow["{\pi_i}"', from=1-1, to=2-1]
                	\arrow["{p_i}", from=1-2, to=2-2]
                	\arrow["{f_i}", from=2-1, to=2-2]
                \end{tikzcd}
                $$p_i \circ f = f_i \circ \pi_i$$
            \end{center}
            Since both $\pi_i$, $f_i$ are continuous so is $f_i \circ \pi_i$.
            Therefore $p_i \circ f$ is also continuous.
        \item Let $J_0$ be the set of $i\in I$ such that $f_i$ is not surjective.
            A topological basis of $X$ can be given by sets of the from
            $$ \bigcap_{j \in J} \pi_{j}^{-1} \left(U_j\right),$$
            where $J$ is a finite subset of $I$ containing $J_0$.
            Note that
            $$ f\left(\bigcap_{j \in J} \pi^{-1}_{j} (U_j)\right) = \bigcap_{j \in J} p^{-1}_{j}\left(f_j(U_j)\right).$$
            Since $f_j$ is open for any $j$, so must be $f$.
            (Checked on the basis of topology.)
    \end{enumerate}
\end{proofenv}

\begin{propositionenv}\label{thm:closure_pi_set}
    Let $I$ be a set and $(X_i)_{i\in I}$ be a family of topological spaces.
    Let $X$ be the product $(X_i)_{i\in I}$ with initial topology.
    For any $i\in I$, let $A_i \subseteq X_i$ and $C_i = \bar{A}_i$ (in $X_i$).
    Put $A = \prod_{i\in I}A_i$ and $C = \prod_{i\in I}C_i$,
    then
    $$ \bar{A} = C \text{ in } X.$$
\end{propositionenv}
\begin{proofenv}
    First, we prove that $C$ is closed in $X$.
    For any $i\in I$, let $\pi_i : X \rightarrow X_i$ be the projection to the $i$-th coordinate.
    We have
    $$ X\setminus C = \bigcup_{i \in I} \pi_i^{-1}\left(X_i \setminus C_i\right).$$
    For the remaining part, take $x = (x_i)_{i \in I}$ to be an element of $C$.
    We claim that for any open $U \ni x$, $U \cap A \neq \varnothing$
    \begin{center}
        

% Pattern Info
 
\tikzset{
pattern size/.store in=\mcSize, 
pattern size = 5pt,
pattern thickness/.store in=\mcThickness, 
pattern thickness = 0.3pt,
pattern radius/.store in=\mcRadius, 
pattern radius = 1pt}
\makeatletter
\pgfutil@ifundefined{pgf@pattern@name@_44loeb5b3}{
\pgfdeclarepatternformonly[\mcThickness,\mcSize]{_44loeb5b3}
{\pgfqpoint{0pt}{0pt}}
{\pgfpoint{\mcSize+\mcThickness}{\mcSize+\mcThickness}}
{\pgfpoint{\mcSize}{\mcSize}}
{
\pgfsetcolor{\tikz@pattern@color}
\pgfsetlinewidth{\mcThickness}
\pgfpathmoveto{\pgfqpoint{0pt}{0pt}}
\pgfpathlineto{\pgfpoint{\mcSize+\mcThickness}{\mcSize+\mcThickness}}
\pgfusepath{stroke}
}}
\makeatother
\tikzset{every picture/.style={line width=0.75pt}} %set default line width to 0.75pt        

\begin{tikzpicture}[x=0.75pt,y=0.75pt,yscale=-1,xscale=1]
%uncomment if require: \path (0,300); %set diagram left start at 0, and has height of 300

%Shape: Regular Polygon [id:dp38747178348966016] 
\draw  [line width=2.25]  (312.77,150.33) .. controls (301.5,150.33) and (290.23,150.12) .. (278.96,150.33) .. controls (271.27,150.47) and (257.48,159.96) .. (251.91,164.63) .. controls (249.41,166.73) and (245.86,167.55) .. (243.8,170.08) .. controls (241.21,173.24) and (239.2,176.83) .. (237.03,180.3) .. controls (233.2,186.42) and (229.44,210.16) .. (237.71,215.71) .. controls (256.61,228.41) and (275.93,223.93) .. (297.22,226.61) .. controls (316.7,229.07) and (334.22,234.2) .. (354.02,232.74) .. controls (370.44,231.53) and (384.51,215.23) .. (386.48,199.37) .. controls (387.22,193.45) and (388.9,188.09) .. (385.13,183.02) .. controls (377.39,172.62) and (363.57,173.37) .. (352.67,170.08) .. controls (344,167.46) and (341.35,164.81) .. (336.44,159.86) .. controls (334.23,157.64) and (322.2,149.76) .. (312.77,150.33) -- cycle ;
%Shape: Circle [id:dp10118816507723016] 
\draw  [line width=2.25]  (270,147) .. controls (270,133.19) and (281.19,122) .. (295,122) .. controls (308.81,122) and (320,133.19) .. (320,147) .. controls (320,160.81) and (308.81,172) .. (295,172) .. controls (281.19,172) and (270,160.81) .. (270,147) -- cycle ;
%Shape: Free Drawing [id:dp7937280453890982] 
\draw  [line width=3] [line join = round][line cap = round] (296.74,148.73) .. controls (296.74,149.24) and (296.74,149.76) .. (296.74,150.27) ;
%Shape: Path Data [id:dp8272736758139742] 
\draw  [pattern=_44loeb5b3,pattern size=6pt,pattern thickness=0.75pt,pattern radius=0pt, pattern color={rgb, 255:red, 0; green, 0; blue, 0}][line width=2.25]  (312.77,150.33) .. controls (314.98,150.19) and (317.33,150.52) .. (319.66,151.13) .. controls (317.69,162.97) and (307.4,172) .. (295,172) .. controls (283.14,172) and (273.2,163.74) .. (270.64,152.65) .. controls (273.71,151.28) and (276.62,150.37) .. (278.96,150.33) .. controls (290.23,150.12) and (301.5,150.33) .. (312.77,150.33) -- cycle ;

% Text Node
\draw (288,185.4) node [anchor=north west][inner sep=0.75pt]    {$A$};
% Text Node
\draw (284,132.4) node [anchor=north west][inner sep=0.75pt]    {$x$};


\end{tikzpicture}
    \end{center}
    Without loss of generality, we can assume that $U = \prod_{i\in I}U_i$ such that for all $i\in I$ that $U_i$ open in $X_i$ and $U_i = X_i$ for all but finitely many $i \in I$
    ( $X = \prod_{i\in I}X_i $)
    Since $C_i = \bar{A}_i$, so $x_i \in U_i \cap A_i \neq \varnothing$ for any $i\in I$ hence
    $$ (x_i)_{i\in I} \in \prod_{i\in I} U_i \cap \prod_{i \in I}A_i = A.$$
    So $C$ is the closure of $A$.
\end{proofenv}

\begin{propositionenv}\label{prop:5.1.6}
    Let $X$ be a topological space and $\sim$ be an equivalence relation on $X$.
    Let $Y$ be the quotient of $X$ by equivalence relation considered with the quotient topology and let $\Gamma$ be the graph of $\sim$.
    \begin{enumerate}
        \item If $Y$ is Hausdorff, then $\Gamma$ is closed in $X \times X$.
        \item The converse is true if we assume that $\pi : X\rightarrow Y$ is open.
    \end{enumerate}
\end{propositionenv}
\begin{proofenv}
    Let
    $$ \Delta_Y \coloneq \{(y,y)\mid y \in Y\} \subseteq Y \times Y$$
    be the diagonal of $Y\times Y$.
    Consider
    $$\begin{array}{rrcl}
        \pi \times \pi : & X \times X & \longrightarrow & Y \times Y\\
        & \left(x_1,x_2\right) & \longmapsto & \left(\pi(x_1),\pi(x_2)\right).
    \end{array}$$
    By definition, $\Gamma = \left(\pi \times \pi\right)^{-1} \left(\Delta_Y\right)$, so if $Y$ is Hausdorff, then $\Gamma$ is closed, since $\pi \times \pi$ is continuous by 1. of \ref{prop:5.1.4}.

    \quad Assume that $\pi$ is open.
    By 2. of \ref{prop:5.1.4}, $\pi \times \pi$ is open.
    If $\Gamma$ is closed, then
    $$ X\times X \setminus \Gamma = \left(\pi \times \pi\right)^{-1} \left(Y \times Y \setminus \Delta_Y\right)$$
    is open.
    The image of this set by $\pi \times \pi$ identifies with $Y \times Y \setminus \Delta_Y$, so is also open.
    Hence $\Delta_Y$ is closed, i.e. $Y$ is Hausdorff.
\end{proofenv}

\begin{definitionenv}
    We say that a topological spaces $(X,\mathscr{T})$ is \textbf{disconnected} is there exists $U,V \subseteq X$ open and non-empty such that $U \cap V = \varnothing$ and $X = U \cup V$.
    If $X$ is not disconnected, then we say that it is \textbf{connected}.

    \quad We say that a subset $A \subseteq X$ is connected if $(A,\mathscr{T}|_A)$ is connected.
    
    \quad We say that $(X,\mathscr{T})$ is \textbf{path connected} if for any $x,y\in X$, there exists a continuous mapping $\gamma : [0,1]\rightarrow X$ such that $\gamma(0) = x$ and $\gamma(1) = y$.
\end{definitionenv}
\begin{propositionenv}\label{thm:connectedSpace_iff_clopen}
    Let $(X,\mathscr{T})$ be the topological space.
    $X$ is connected is and only if there does not exist a subset of $X$ different from both $X$ and $\varnothing$ that is simultaneously open and closed.
\end{propositionenv}
\begin{proofenv}
    Assume that there exists $U \subsetneq X$ such that $U \neq \varnothing$ and $U$ is simultaneously open and closed.
    Put $V = X\setminus U$, then $V \cap U = \varnothing$ and $V\cup U =X$, thus $X$ is disconnected, contradiction!
\end{proofenv}

\begin{exampleenv}
    $\sin\left(\frac{1}{x}\right) \cup \{(0,y): -1\le y\le 1\}$
 is connected but not path connected.
\end{exampleenv}

\begin{propositionenv}\label{thm:IsPathConnected.isConnected}
    Let $(X,\mathscr{T})$ be a topological space.
    If $X$ is path connected, then $X$ is connected.
\end{propositionenv}
\begin{proofenv}
    We prove by contradiction.
    We write $X$ as the union of two disjoint non-empty open subsets ($X = U\cup V,\ U\cap V =\varnothing,\ U,V \neq \varnothing$) that $X$ is disconnected but path connected.
    Pick $x\in U, y \in V$ and $\gamma:[0,1]\rightarrow X$ such that $\gamma(0) = x$ and $\gamma(1) = y$.
    Then $\gamma^{-1}(U)$, $\gamma^{-1}(V)$ are disjoint open subsets of $[0,1]$ containing $0$ and $1$ respectively and $[0,1] = \gamma^{-1}(U)\cup \gamma^{-1}(V)$.
    Let
    $$ a = \sup \{t \in [0,1] \mid \interval[open right]{0}{t} \subseteq \gamma^{-1}(U)\}.$$
    Since $\gamma^{-1}(U)$ is closed, $a \in \gamma^{-1}(U)$.
    Since $\gamma^{-1}(U)$ is open, there exists $\varepsilon > 0$ such that $\interval[open]{a-\varepsilon}{a+\varepsilon} \subseteq \gamma^{-1}(U)$, contradiction!
\end{proofenv}

\begin{remark}
    Note that we proved that all intervals are connected.
\end{remark}

\begin{definitionenv}
    Let $(X,\mathscr{T})$ be a topological spaces.
    We call \textbf{connected component} of $X$ any non-empty maximal element of the set of all connected subsets of $X$ (equipped with the order of inclusion of sets).
    An empty topological space is connected but has no connected component.
\end{definitionenv}

\begin{lemmaenv}\label{thm:IsConnected.biUnion_of_reflTransGen}
    Let $(X,\mathscr{T})$ be a topological space, $\left(A_i\right)_{i\in I}$ family of non-empty connected subsets of $X$.
    If there exists $i_0 \in I$ such that for all $i \in I$, $A_{i} \cap A_{i_0} \neq \varnothing$, then $\bigcup_{i\in I}A_i$ is connected.
\end{lemmaenv}
\begin{proofenv}
    Suppose $U,V \subseteq X$, $U \cup V = X$, $U \cap V \neq \varnothing$,
    $$ \bigcup_{i\in I} A_i \subseteq U \cup V, \bigcup_{i\in I}A_i \cap \left(U \cap V\right) = \varnothing.$$
    $A_{i_0}$ is connected so $A_{i_0} \subseteq U$ or $A_{i_0} \subseteq V$.
    Without loss of generality, assume that $A_{i_0} \subseteq U$.
    For any $i \in I$, $\varnothing \neq A_i\cap A_{i_0} \subseteq A_{i_0} \subseteq U$.
    So for any $i \in I$, $A_i\cap U \neq \varnothing$, $A_i$ is connected $A_i \subseteq U$.
    Thus 
    $$ \bigcup_{i\in I}A_i \subseteq U.$$
\end{proofenv}

\begin{propositionenv}
    Let $(X,\mathscr{T})$ be a topological space.
    Then the following binary relation $\sim$ is an equivalence relation:
    \begin{center}
        $x \sim y \ \Leftrightarrow$ there exists a connected subset $A$ of $X$ such that $\{x,y\} \subseteq A$. 
    \end{center}
    Moreover, the equivalence classes of this relation are exactly the connected components of $(X,\mathscr{T})$.
    
    \quad In particular, any $x \in X$ belongs to a \textit{unique} connected component (called the \textbf{connected component} of $X$).
\end{propositionenv}
\begin{proofenv}
    $\sim$ is already reflexive and symmetric.
    We prove that $\sim$ is transitive.
    Let $x,y, z$ be elements of $X$ such that $\{x,y\}\subseteq A$ and $\{y,z\} \subseteq B$.
    One has $A \cap B \neq \varnothing$ case $y \in A \cap B$.
    So by lemma \ref{thm:IsConnected.biUnion_of_reflTransGen}, $A \cup B$ is connected and since ${x,z} \subseteq A \cup B$, one has $x \sim z$.

    \quad Let $C$ be a connected component of $X$ and $x \in X$.
    By definition, $x\sim y$ for any $y \in C$.
    If $z \in X$ is such that $x \sim Z$, then there exists connected $A \subseteq X$ such that $\{x,z\} \subseteq A$ and $A \cup C$ is connected.
    Since by definition, $C$ is maximal, we conclude that $A \subseteq C$ and $z \in C$, so $C$ is equal to the equivalence class of $x$.

    \quad Conversely, let $x \in X$ and $C$ be the equivalence class.
    By definition, any connected subset of $X$ containing $x$ must be contained in $C$.
    
    \quad It remains to check that $C$ is connected.
    
    \quad For any $y \in c$, there exists a connected set $A_y \supseteq \{x,y\}$.
    By lemma \ref{thm:IsConnected.biUnion_of_reflTransGen}, $\displaystyle \bigcup_{y \in C} A_y$ is connected and contains $C$.
    We conclude that
    $$C = \bigcup_{y \in C} A_y.$$
    Hence $C$ is connected.
\end{proofenv}

\begin{propositionenv}\label{prop:5.1.15}
    Let $(X,\mathscr{T})$ be a topological space.
    If $A \subseteq X$ is a connected subset of $X$, then $\bar{A}$ is also connected.
    In particular, any connected component of $X$ is closed.
\end{propositionenv}
\begin{proofenv}
    Suppose, by contradiction, that $U,V \subseteq X$ are open such that
    $$ (U \cap V) \cap \bar{A} = \varnothing,\ \bar{A} \subseteq U \cup V,\ U \cap \bar{A} \neq \varnothing,\ V \cap \bar{A} \neq \varnothing.$$
    But in this situation, one would also have $U \cap A \neq \varnothing$, $V \cap A \neq \varnothing$, $(U \cap V) \cap A =\varnothing$, $A \subseteq U \cup V$ which would contradict the assumption that $A$ is connected.
\end{proofenv}


\section{Topological Groups}

\begin{definitionenv}
    Let $(G,*)$ be a group with a topology $\mathscr{T}$.
    If the mapping
    $$\begin{array}{rrcl}
        \psi : & G\times G & \longrightarrow & G\\
        & (x,y) & \longmapsto & x*\iota(y)
    \end{array}$$
    is continuous, then we say that $(G,*,\mathscr{T})$ is a \textbf{topological group}.
\end{definitionenv}

\begin{definitionenv}
    Let $(G,*)$ be a group and $X$ be a set.
    Recall that a mapping $\varphi: G \times X \rightarrow X$ is called a left action (resp. right action) if the following conditions are satisfied
    \begin{enumerate}
        \item For any $x \in X$, $\varphi(e,x) = x$ ($e$ is the neutral element of $G$).
        \item For any $(a,b) \in G\times G$ and any $x \in X$,
            $$ \varphi(a*b, x) = \varphi(a,\varphi(b,x))\ (\text{resp. } \varphi(a*b, x) = \varphi(b,\varphi(a,x))).$$
    \end{enumerate}
    For any $a \in G$, we denote by $L_a : X \rightarrow X$ (resp. $R_a : X \rightarrow X$) the mapping sends any $x \in X$ to $\varphi(a,x)$.

    \quad If $(G,*,\mathscr{T})$ is a topological space, then we say that a left action (resp. right action) is continuous if $\varphi : G \times X \rightarrow X$ is continuous as a mapping when we consider $G \times X$ with the product topology.
\end{definitionenv}

\begin{propositionenv}\label{prop:5.2.3}
    Let $(G,*,\mathscr{T})$ be a topological group and $X$ be a topological space.
    Let $\varphi: G \times X \rightarrow X$ be a continuous left action (resp. right action) on $X$.
    Then for any $a \in G$, the mapping $L_a$ (resp. $R_a$) is a homeomorphism.
\end{propositionenv}
\begin{proofenv}
    We prove the left action case (right action case is similar).
    The mapping $L_a : X\rightarrow X$, $x \mapsto \varphi(a,x)$ can be written as the components
    $$ X \overset{i_a}{\longrightarrow} G \times X \overset{\varphi}{\longrightarrow} X$$
    where $i_a : X \rightarrow G\times X$ sends $x$ to $(a,x)$.
    Note that $i_a$ composed with the projections $G\times X \rightarrow G$ and $G \times X \rightarrow X$ give us a constant mapping and the identity on $X$ which are both continuous.
    Since projections are always continuous and open, we deduce that $i_a$ is continuous and hence also $L_a = i_a \circ \varphi$ is continuous.

    \quad By definition, 
    $$ L_a \circ L_{\iota(a)} = L_{\iota(a)} \circ L_{a} = \id_{X},$$
    we deduce that $L_a$ is a homeomorphism, cause $L_{\iota(a)}$ is continuous by the first half of the proof.
\end{proofenv}

\begin{propositionenv}
    Let $(G,*,\mathscr{T})$ be a topological group.
    \begin{enumerate}
        \item The inverse mapping $\iota : G\rightarrow G$ is a homeomorphism.
        \item $G \times G \rightarrow G$, $(a,b) \mapsto a*b$ is continuous and in particular, it defines a left action on $G$ in itself.
        \item $G \times G \rightarrow G$, $(a,b) \mapsto b*a$ is continuous and in particular, it defines a right action on $G$ in itself.
        \item For any $g \in G$,
            $$ L_g : G \rightarrow G, \ x \mapsto g*x,$$
            $$ R_g : G \rightarrow G, \ x \mapsto x*g,$$
            are homeomorphisms on $G$.
    \end{enumerate}
\end{propositionenv}
\begin{proofenv}
    Let $e \in G$ be the neutral element,
    and $p_1 : G \times G \rightarrow G$ and $p_2 : G \times G \rightarrow G$ be the two projections,
    namely $$ p_1(a,b) = a,\ p_2(a,b) = b.$$
    We notice that $p_1$ and $p_2$ are both continuous and open.
    \begin{enumerate}
        \item The mapping $\lambda : G \rightarrow G \times G$, $x \mapsto (e,x)$ is continuous since its composite with each of the projections $p_1$ and $p_2$ is continuous.
            Therefore, the mapping $\iota : G \rightarrow G$, which is the composite 
            $$\begin{array}{ccccl}
                G & \overset{\lambda}{\longrightarrow} & G \times G & \overset{\varphi}{\longrightarrow} & G\\
                x & \longmapsto & (e,x) & \longmapsto & e* \iota(x)
            \end{array}$$
            is continuous.
            Since $\iota \circ \iota = \id_{G}$, we deduce that $\iota$ is a homeomorphism of $G$.
        \item Let $\tau : G \times G \rightarrow G \times G$, $(x,y) \mapsto (x, \iota(y))$ one has that $p_1 \circ \tau = p_1$, $p_2 \circ \tau = \iota \circ p_2$,
            and therefore $\tau$ is continuous.
            The mapping
            $$ \begin{array}{rrcl}
                G \times G & \longrightarrow & G\\
                (x,y) & \longmapsto & x*y
            \end{array}$$
            equals $\varphi \circ \tau$ so $(x,y) \mapsto x*y$ is continuous.
        \item $G \times G \rightarrow G \times G$, $(x,y) \mapsto (y,x)$ is continuous,
            since its composition with both $p_1$ and $p_2$ are continuous.
            The mapping $G \times G \rightarrow G$, $(x,y) \mapsto y*x$ is the composite of $(x,y) \mapsto (y,x)$ and $(x,y)\mapsto x*y$.
            Hence it is continuous.
        \item Follows from the proposition \ref{prop:5.2.3}.
    \end{enumerate}
\end{proofenv}

\begin{remark}
    Let $(G,*)$ be a group with topology $\mathscr{T}$.
    We can prove that $(G,*,\mathscr{T})$ is a topological group if and only if 
    $$\begin{array}{rrclcrrcl}
        \iota : & G & \longrightarrow & G & \quad & * : & G \times G & \longrightarrow & G \\
        & x & \longmapsto & \iota(x) & \quad & & (x,y) & \longmapsto & x*y
    \end{array}$$
    are both continuous.
\end{remark}

\begin{exampleenv}
    Let $(K, \left|\ \cdot\ \right|)$ be a valued field.
    The additive group $(K,+)$ and the multiplicative group $(K^\times,\cdot)$ are both topological groups.
\end{exampleenv}

\begin{exampleenv}
    Let $(K, \left|\ \cdot\ \right|)$ be a valued field, and $(V,\|\cdot\|)$ be a normed vector space over $K$.
    The additive group $(V,+)$ forms a topological group.
    Moreover, the mapping
    $$ \begin{array}{rrcl}
        K^\times \times V & \longrightarrow & V\\
        (a,x) & \longmapsto & ax
    \end{array}$$
    is a continuous left action.
\end{exampleenv}

\begin{exampleenv}
    Let $(K, \left|\ \cdot \ \right|)$ be a complete valued field,
    and $(V,\|\cdot\|)$ be a Banach space over $K$,
    $\mathcal{L}(V,V)$ vector space of bounded $K$-linear endomorphisms.
    $\mathcal{L}(V,V)$ is equipped with the operator norm.
    The bilinear mapping
    $$ \begin{array}{rcl}
        \mathcal{L}(V,V) \times \mathcal{L}(V,V) & \longrightarrow & \mathcal{L}(V,V)\\
        (\varphi,\psi) & \longmapsto & \varphi \circ \psi
    \end{array}$$
    is bounded and hence continuous.
    Let $\GL(V)$ be subset of $\mathcal{L}(V,V)$ of bounded $K$-linear isomorphisms.

    \quad Note that by Banach's open mapping theorem, for any $f \in \GL(V)$, then $f^{-1} \in \GL(V)$.
    Hence $\GL(V)$ forms a group under composition of mappings.
    It also has been shown that the mapping of inversion $\iota : \GL(V) \rightarrow \GL(V)$ is of class $C^\infty$.
    So it is continuous.
    Therefore $\GL(V)$ is a topological group.
    The mapping
    $$ \begin{array}{rcl}
        \GL(V) \times V & \longrightarrow & V\\
        (f,x) & \longmapsto & f(x)
    \end{array}$$
    is a continuous left action.
    In fact the mapping
    $$ \begin{array}{rcl}
       \mathcal{L}(V, V) \times V & \longrightarrow & V\\
        (f,x) & \longmapsto & f (x)
    \end{array}$$
    is a bounded linear mapping (so continuous).
\end{exampleenv}

\begin{propositionenv}\label{prop:5.2.9}
    Let $(G,*,\mathscr{T})$ be a topological group,
    $X$ be a topological space,
    $\varphi : G \times X \rightarrow X$ be a continuous left action.
    Let $U \subseteq G$ be a subset, $V \subseteq X$ be open,
    then
    $$\{\varphi(a, x) : (a, x) \in U \times V\}$$
    is open in $X$.
\end{propositionenv}
\begin{proofenv}
    One has 
    $$ \{\varphi(a, x) : (a, x) \in U \times V\} = \bigcup_{a \in U} L_a(V).$$
    $L_a (V)$ is open since $L_a$ is a homeomorphism.
\end{proofenv}

\begin{propositionenv}\label{prop:5.2.10}
    Let $(G,*,\mathscr{T})$ be a topological group,
    $e \in G$ be the neutral element,
    $e \in U \subseteq G$ open.
    \begin{enumerate}
        \item There exists an open $V \subseteq G$ such that $e \in V$, $\iota (V) = V$, and
            $$V * V = \{x * y \mid (x,y) \in V \times V\} \subseteq U.$$
        \item For any element $a \in U$, there exists $e \in V \subseteq U$ open such that 
            $$ a * V =\{a * x \mid x \in V \} \subseteq U,\ V * a = \{x * a \mid x \in V\} \subseteq U.$$
        \item For any $b \in U$, there exists $e \in V \subseteq U$ open such that
            $$ b * V * \iota(b) = \{b * x * \iota(b) \mid  x \in V\} \subseteq U.$$
    \end{enumerate}
\end{propositionenv}
\begin{proofenv}
    \quad 
    \begin{enumerate}
        \item $\varphi : G \times G \rightarrow G$, $(x,y) \mapsto x* \iota(y)$ is continuous and it sends $(e,e)$ to $e$.
            Hence $\varphi^{-1}(U) \ni (e,e)$ is open.
            There exists $V_1,V_2 \subseteq G$ open neighborhoods of $e$ such that
            $$ V_1 \times V_2 \subseteq \varphi^{-1}(U).$$
            We put $V = (V_1\cap V_2) \cap \iota\left(V_1 \cap V_2\right)$.
            It is an open neighborhood of $e$ such that $\iota(V) = V$ and $V \times V \subseteq \varphi^{-1}(U)$ so $V * V = V * \iota(V) \subseteq U$.
        \item We take 
            $$ V = L_{\iota(a)} (U) \cap R_{\iota(a)} (U) \cap U$$
            open neighborhood of $e$.
            Moreover,
            $$ a * V = L_{a} (V) \subseteq U,\quad V * a = R_{a} (V) \subseteq U.$$
        \item Put 
            $$ V = \left(\iota(b) * U * b\right) \cap U = R_{b} \left(L_{\iota(a)} (U)\right) \cap U.$$
            It is an open neighborhood of $e$ contained in $U$ ($\iota(b) * e * b = e$).
            Moreover,
            $$ b * V * \iota(b) \subseteq U.$$
    \end{enumerate}
\end{proofenv}

\begin{propositionenv}
    Let $(G,*,\mathscr{T})$ be a topological group,
    $e \in G$ be the neutral element.
    The topological group $(G,*,\mathscr{T})$ is Hausdorff if and only if the intersection of all open neighborhoods of $e$ is equal to $\{e\}$.
\end{propositionenv}
\begin{proofenv}
    Suppose that $(G,\mathscr{T})$ is Hausdorff.
    Let $x \in G$, $x \neq e$.
    There exists $e \in U_x$ open and $x \notin U_x$.
    Hence,
    $$ \bigcap_{\substack{x \in G\\ x \neq e}} U_x = \{e\}.$$

    \quad Assume that the intersection of all open neighborhoods of $e$ is equal to $\{e\}$.
    Let $x, y\in G$ such that $x \neq y$, then $x * \iota(y) \neq e$.
    There exists an open $U \ni e$ such that $x * \iota(y) \notin U$.
    By \ref{prop:5.2.10}, there exists $V$, such that $V * \iota(V) \subseteq U$.
    We claim that $L_{x} (V) \cap L_{y}(U) = \varnothing$.
    Indeed, if $(a,b) \in V$ such that $x * a = y * b$, then
    $$a * \iota(b) = \iota(x) * y \in U,$$
    contradiction.
\end{proofenv}

\begin{propositionenv}\label{prop:5.2.12}
    Let $(G,*,\mathscr{T})$ be a topological group,
    $H$ be a subgroup of $G$.
    \begin{enumerate}
        \item $\bar{H}$ is also a subgroup of $G$.
        \item If $H$ is normal, then $\bar{H}$ is also normal.
        \item If $H$ is open, then $H$ is also closed.
    \end{enumerate}
\end{propositionenv}
\begin{proofenv}
    \quad
    \begin{enumerate}
        \item 
            $$\begin{array}{rrcl}
                \varphi : & G \times G & \longrightarrow & G \\
                & (x,y) & \longmapsto & x * \iota(y)
            \end{array}$$
            is a continuous mapping, so $\varphi^{-1}(\bar{H})$ is closed in $G \times G$.
            Also, 
            $$ \varphi^{-1}(\bar{H}) \supseteq \varphi^{-1}(H) \supseteq H \times H.$$
            By proposition \ref{thm:closure_pi_set}, $\varphi^{-1}(\bar{H}) \supseteq \bar{H} \times \bar{H}$.
            So $\bar{H}$ is also a subgroup of $G$.
        \item Suppose $H$ is normal, namely, for any $g \in G$,
            $$H \subseteq g * H * \iota(g).$$
            Note that this condition is equivalent to
            $$\forall g \in G,\ g * H * \iota(g) = H.$$
            Let now $\psi_g : G \rightarrow G$, $x \mapsto g * x* \iota(g)$.
            Note that $\psi_{g} = R_{\iota(g)} \circ L_g$, so $\psi_{g}$ is continuous.
            Since $H$ is normal,
            $$ H \subseteq \psi^{-1}_{g} (H) \subseteq \psi^{-1}_{g} (\bar{H}).$$
            So $\bar{H} \subseteq \psi^{-1}_{g} (\bar{H})$, so $\bar{H}$ is normal.
        \item If $H$ is open, then $g * H$ is open for any $g \in G$ and
            $$ H = G \setminus \bigcup_{\substack{g \in G\\g \notin H}} g * H.$$
            So $H$ is closed.
    \end{enumerate}
\end{proofenv}

\begin{definitionenv}
    Let $G$ be a group, $X$ be a set, $\varphi: G \times X \rightarrow X$ be a left action of $G$ on $X$ (resp. right action).
    \begin{enumerate}
        \item If for any $g \in G$, $e \neq g$, there exists $x \in X$ such that $\varphi(g,x) \neq x$,
            then we say that $\varphi$ is \textbf{faithful}.
        \item If for any $(x,y) \in X$, there exists $g \in G$ such that $\varphi(g,x) = y$,
            then we say that $\varphi$ is \textbf{transitive},
            and that $X$ is a \textbf{homogeneous space} of $G$.
        \item For any $x \in X$, we denote by
            $$ \stab (x) = \{g \in G \mid \varphi(g,x) = x\}.$$
            This is a subgroup of $G$ called the \textbf{stabilizer} of $x$ with respect to $\varphi$.
    \end{enumerate}
\end{definitionenv}

\begin{exampleenv}
    Let $(G,*,\mathscr{T})$ be a topological group,
    $H$ be a subgroup of $G$,
    $G/H$ the set of left cosets of $G$ with respect to $H$.
    Let $\pi : G \rightarrow G/H$ be a projection.
    We equipped $G/H$ with the quotient topology,
    i.e. $V \subseteq G/H$ is open if and only if $\pi^{-1} (V)$ is open in $G$.

    \quad Note that $\pi : G \rightarrow G/H$ is open.
    In fact if $U \subseteq G$ open, then
    $$ \pi^{-1} \left(\pi(U)\right) = \bigcup_{h \in H} U * h$$
    is an open subset of $G$.
    The mapping
    $$\begin{array}{rrcl}
        \psi : & G \times (G/H) & \longrightarrow & G/H\\
        & (g, x * H) & \longmapsto & (g * x) * H
    \end{array}$$
    defines a left action of $G$ on $G/H$.
    It turns out that this action is continuous.
    Indeed, if $V \subseteq G/H$ is open, then
    $$\begin{aligned}
        \psi^{-1} (V) = & \{(g, x * H) \in G \times G/H \mid (g * x) * H \in V\}\\
        = & \{(g, x * H) \in G \times G/H \mid g * x \in \pi^{-1} (V)\}.
    \end{aligned}$$
    Let
    $$\begin{array}{rrcl}
        m : & G \times G & \longrightarrow & G\\
        & (x,y) & \longmapsto & x * y
    \end{array}$$
    It is continuous.
    Hence $m^{-1} (\pi^{-1} (V))$ is open in $G \times G$.

    \quad Note that $\psi^{-1} (V)$ is the image of $m^{-1}(\pi^{-1} (V))$ by
    $$\begin{array}{rcl}
        G \times G & \longrightarrow & G/H\\
        (x,y) & \longmapsto & x * y * H
    \end{array}$$
    is open.
    Thus $\psi^{-1} (V)$ is an open subset of $G \times G/H$.
    By definition, $\psi$ is transitive, hence $G/H$ is a homogeneous space of $G$.
\end{exampleenv}

\begin{exampleenv}
    Let $(G,*,\mathscr{T})$ be a topological group,
    $H$ be a subgroup of $G$.
    $$\begin{array}{rrcl}
        \varphi : & G \times G & \longrightarrow & G\\
        & (x,y) & \longmapsto & x * \iota(y)
    \end{array}$$
    is continuous.
    We equipped the group $G/H$ with the quotient topology.
    Let
    $$\begin{array}{rrcl}
        \psi : & (G/H) \times (G/H) & \longrightarrow & G/H\\
        & (x * H, y * H) & \longmapsto & (x * \iota(y)) * H.
    \end{array}$$
    We claim that $\psi$ is continuous.
    Let $\pi : G \rightarrow G/H$ be the projection (continuous and open).
    If $V \subseteq G/H$ is open, then $\psi^{-1}$ is open, then $\psi^{-1}(V)$ is the image of $\varphi^{-1}(\pi^{-1}(V))$ by the open mapping.
    $$\begin{array}{rcl}
        G \times G & \longrightarrow & (G/H, G/H)\\
        (x, y) & \longmapsto & (x * H, y * H).
    \end{array}$$
    Hence it is open in $(G/H) \times (G/H)$.
\end{exampleenv}

\begin{exampleenv}
    Let $(V, \langle\ ,\ \rangle)$ be the finite-dimensional inner product space over $\RR$,
    $$ S =\{x \in V \mid \langle x, x \rangle = 1\}$$
    be the unit sphere of $(V,\langle\ ,\ \rangle)$.
    Let $O(V,\langle\ ,\ \rangle)$ be the orthogonal group of $(V,\langle\ ,\ \rangle)$ ($\RR$-linear isomorphisms that are $\RR$-linear isometries).
    Then $O(V,\langle\ ,\ \rangle)$ acts continuously on $S$.
    We choose two orthogonal bases $\left(e_i\right)_{i=1}^n$, $\left(e_i'\right)_{i = 1}^{n}$ of $(V,\langle\ ,\ \rangle)$ such that $x = e_1$ and $y = e_1'$.
    Then the mapping
    $$\begin{array}{rrcl}
        f : & V & \longrightarrow & V\\
        & \lambda_1 e_1 + \cdots + \lambda_n e_n & \longmapsto & \lambda_1 e_1' + \cdots + \lambda_n e_n'
    \end{array}$$
    is a linear isometry that takes $x$ to $y$.
    So $S$ is a homogeneous space of the group $O(V,\langle\ ,\ \rangle)$.
\end{exampleenv}

\begin{propositionenv}\label{prop:5.2.17}
    Let $(G,*,\mathscr{T})$ be a topological group,
    $X$ be a topological space,
    $\varphi : G \times X \rightarrow X$ a continuous left action and $\varphi$ transitive.
    Let $x_0 \in X$, and $H =\stab (x_0)$.
    Then the mapping
    $$\begin{array}{rrcl}
        f : & G/H & \longrightarrow & X\\
        & a * H & \longmapsto & \varphi (a,x_0)
    \end{array}$$
    is bijective and continuous.
\end{propositionenv}
\begin{proofenv}
    Let $\pi : G \rightarrow G/H$ be the projection, $g: G\rightarrow X$, $a \mapsto \varphi (a,x_0)$.
    Since $\varphi$ and the mapping
    $$\begin{array}{rcl}
        G & \longrightarrow & G \times X\\
        a & \longmapsto & (a,x_0)
    \end{array}$$
    is continuous, we get that $g$ is continuous.
    Moreover,
    $$ g = f \circ \pi.$$
    So for any open set $U \subseteq X$, we get that
    $$ g^{-1} (U) = \pi^{-1} \left(f^{-1} (U)\right)$$
    is open.
    Hence $f^{-1}(U)$ must be open too (in $G/H$).
    It remains to check that $f$ is a bijection.
    Since $g$ is a surjective mapping (by transitivity of $\varphi$), $f$ is injective.
    If $a,b \in G$, are such that $\varphi (a,x_0) = \varphi (b,x_0)$, then $\varphi (\iota(a) * b, x_0) = x_0$.
    Namely, $\iota(a) * b \in H$, thus $a * H = b * H$.
    So $f$ is also injective.
\end{proofenv}

\begin{propositionenv}
    Let $(G,*,\mathscr{T})$ be a topological group,
    $X$ be a topological space,
    $\varphi : G \times X \rightarrow X$ a continuous left action.
    If $X$ is Hausdorff, then for any $x_0 \in X$, $\stab(x_0)$ is a closed subgroup of $G$.
\end{propositionenv}
\begin{proofenv}
    $$\begin{array}{rrcl}
        f : & G & \longrightarrow & X\\
        & a & \longmapsto & \varphi (a,x_0)
    \end{array}$$
    is continuous.
    We have 
    $$ f^{-1} (\{x_0\}) = \stab(x_0).$$
    ($\{x_0\}$ is closed since $X$ is Hausdorff.)
\end{proofenv}

\begin{propositionenv}
    Let $(G,*,\mathscr{T})$ be a topological group,
    $H$ be a subgroup of $G$,
    $\pi : G \rightarrow G/H$ be the projection.
    Then $G/H$ with the quotient topology is Hausdorff if and only if $H$ is closed.
\end{propositionenv}
\begin{proofenv}
    Let $e \in G$ be the neutral element.
    We suppose $G/H$ is Hausdorff,
    then $\{e * H\}$ is closed in $G/H$.
    Hence $H = \pi^{-1} (e * H)$ is closed.
    
    \quad Conversely, suppose that $H$ is closed.
    Let $\Gamma$ be the graph of the following equivalence relation:
    $$ a \sim b \Leftrightarrow \iota(a) * b \in H.$$
    Let $(a,b) \in G \times G$ such that $\iota(a) * b \notin H$.
    The mapping
    $$\begin{array}{rrcl}
        f : & G \times G & \longrightarrow & G\\
        & (x,y) & \longmapsto & \iota(x) * y
    \end{array}$$
    is continuous.
    There exists $U_a$, $U_b$ open neighborhoods of $a,b$ respectively such that
    $$ f (U_a \times U_b) \subseteq G/H.$$
    Thus
    $$ (a,b) \in U_a \times U_b \subseteq (G \times G)/\Gamma.$$
    By proposition \ref{prop:5.1.6}, we obtain that $G/H$ is Hausdorff.
\end{proofenv}

\begin{definitionenv}
    Let $(G,*,\mathscr{T})$, $(G',*',\mathscr{T}')$ be topological groups.
    If a homomorphism of groups $f : G \rightarrow G'$ is continuous,
    then we say that $f$ is a \textbf{homomorphism of topological groups}.
    If $f$ is an isomorphism of groups that is also a homeomorphism,
    then we say that $f$ is an \textbf{isomorphism of topological groups}.
\end{definitionenv}

\begin{propositionenv}
    Let $(G,*,\mathscr{T})$, $(G',*,\mathscr{T}')$ be topological groups,
    and let $f : G \rightarrow G'$ be a homomorphism of topological groups.
    Assume that $f$ is surjective and open.
    Let $N$ be the kernel of $f$,
    and $\bar{f} : G/N \rightarrow G$ be the isomorphism of groups sending $x$ to $f(x)$.
    Then $\bar{f}$ is an isomorphism of topological groups.
    \begin{center}
        \begin{tikzcd}
        	G & \\
        	{G/\ker(f)} & {G'}
        	\arrow["\pi"', from=1-1, to=2-1]
        	\arrow["f", from=1-1, to=2-2]
        	\arrow["{\bar{f}}"', dashed, from=2-1, to=2-2]
        \end{tikzcd}
    \end{center}
\end{propositionenv}
\begin{proofenv}
    Let $\pi : G \rightarrow G/N$ be te projection.
    We know that $\bar{f}$ is an isomorphism of groups.
    It remains to check that $\bar{f}$ is a homeomorphism.

    \quad For any open $V \subseteq G/N$, one has $\bar{f} (V) = f(\pi^{-1}(V))$ is open in $G'$.
    This means that $\bar{f}$ is an open mapping.
    If $V'$ is open in $G'$,
    one has
    $$ \pi\left(f^{-1} (V')\right) = f^{-1} (V')$$
    open in $G$.
    Hence $\bar{f}^{-1}(V')$ is open in $G/N$,
    so $\bar{f}$ is continuous.
\end{proofenv}

\begin{propositionenv}
    Let $(G,*,\mathscr{T})$ be a topological group, where $G$ is connected.
    Let $e$ be the neutral element of $G$, $U \ni e$ open in $G$.
    Then $G$ is generated by $U$.
    Namely, any element of $G$ can be written as the composition of a family of elements of $U$.
\end{propositionenv}
\begin{proofenv}
    By replacing $U$ by $U \cap \iota(U)$, we may assume that $U = \iota(U)$.
    For any positive integer $n$, let $U_n = \underset{n\text{-times}}{\underbrace{U*\cdots*U}}$.
    By \ref{prop:5.2.9} we obtain (By induction on $n$) that $U_n$ is an open subset of $G$.
    Let 
    $$H = \bigcup_{n\in \NN_{\ge 1}} U_n$$
    be a sub-monoid of $G$.
    Since $U = \iota(U)$, $H$ is also a subgroup of $G$.
    Hence $H$ is an open subgroup of $G$, which is also closed by \ref{prop:5.2.12}.
    Since $G$ is connected, $H = G$.
\end{proofenv}

\begin{propositionenv}
    Let $(G,*,\mathscr{T})$ be a topological group,
    $H$ be a subgroup of $G$.
    If $H$ and $G/H$ are both connected, so is $G$.
\end{propositionenv}
\begin{proofenv}
    Suppose, by contradiction, that $U_1$ and $U_2$ are open such that $U_1\cap U_2 = \varnothing$, $U_1 \cup U_2 = G$.
    Let $\pi : G \rightarrow G/H$ be the projection.
    $\pi$ is open.
    $$G/H = \pi(U_1) \cup \pi(U_2),$$
    but since $G/H$ is connected, $\pi(U_1) \cap \pi(U_2) \neq \varnothing.$
    Let $a_1$ and $a_2$ be in $U_1$ and $U_2$ respectively such that $\pi(a_1) = \pi(a_2)$.
    Let $H' = a_1 * H = a_2 * H$.
    $H$ is connected, so is $H'$.
    But $H' \cap U_1 \neq \varnothing$, $H' \cap U_2 \neq \varnothing$, $H' \cap U_1 \cup H' \cap U_2 = H'$, contradiction.
\end{proofenv}

\begin{propositionenv}
    Let $(G,*,\mathscr{T})$ be a topological group,
    and $H$ connected component of the neutral element $e \in G$.
    Then $H$ is a closed and normal subgroup of $G$,
    and for any $x \in G$, $x * H$ is the connected component of $x$.
\end{propositionenv}
\begin{proofenv}
    By \ref{prop:5.1.15}, $H$ is closed.
    We check that $H$ is a subgroup of $G$.
    Let $a \in H$, $\iota(a) * H$ is connected since $L_{\iota(a)}$ is a homeomorphism.
    Since $e \in \iota(a) * H$, so $H = \iota(a) * H$, $H$ is indeed a subgroup.
    
    \quad We want to check that $H$ is normal.
    For any $g \in G$, $g * H * \iota(g)$ is a connected subset of $G$ containing $e$.
    So $g * H * \iota(g) \subseteq H$, i.e. $H$ is normal.
    For any $x \in G$, the set $x * H$ is a connected subset of $G$ containing $x$.
    Conversely, if $A$ is connected and $x \in A$, then $\iota(x) * A$ is connected containing $e$,
    so $\iota(x) * A \subseteq H$ and so
    $$ A = x * \iota (x) * A \subseteq x * H.$$
\end{proofenv}

\begin{propositionenv}\label{prop:5.2.25}
    Let $f : \RR \rightarrow \RR$ be a continuous homomorphism of the additive group of the real numbers, then $f$ is of the form
    $$ f (x) = x f(1).$$
\end{propositionenv}
\begin{proofenv}
    By the universal property of $(\ZZ, +)$, one has:
    $$ \forall n \in \ZZ,\ f(n) = n f(1).$$
    For any $m \in \ZZ_{>0}$, $n f(1) = f(n) = f(nm/m) = m f(n/m)$.
    So
    $$ f(n/m) = n/m f(1).$$
    Since $f$ is continuous and $\QQ$ is dense in $\RR$, we prove the result.
\end{proofenv}

\begin{corollaryenv}\label{cor:5.2.26}
    Let $g : \RR^\times \rightarrow \RR_{\ge 0}$ be continuous, homomorphism of multiplicative groups.
    There exists $t \in \RR$ such that
    $$g(y) = |y|^t,\ \forall y \in \RR.$$
\end{corollaryenv}
\begin{proofenv}
    Let $f : \RR \rightarrow \RR$, $x \mapsto \ln\left(g(e^x)\right)$.
    $$ f(x_1 + x_2) = \ln(g(e^{x_1 + x_2})) = \ln (g(e^{x_1})g(e^{x_2})) = f(x_1) + f(x_2).$$
    So $f$ is continuous homomorphism, by \ref{prop:5.2.25}
    $$f(x) = x f(1) = x \ln g(e),$$
    so for any $y \in \RR_{\ge 0}$, $g(y) = g(e^{\ln y}) = e^{\ln g(e) \ln y}.$
    $g(1) = g(-1) g(-1) = 1$, $g(-1) > 0$ so $g(-1) = -1$.
    So for any $y \in \RR^\times$, $g(y) = |y|^{\ln g(e)}$.
\end{proofenv}


\section{Locally Compact Hausdorff Spaces}

\begin{definitionenv}
    Let $(X,\mathscr{T})$ be a topological space.
    We say that $(X,\mathscr{T})$ is \textbf{locally compact} if for any $x \in X$,
    there exists a compact neighborhood $C_x$.
\end{definitionenv}

\begin{exampleenv}
    The real line $\RR$ (with standard topology induced $|\cdot|$).
\end{exampleenv}

\begin{definitionenv}
    Let $(X,\mathscr{T})$ be a topological space,
    $f : X \rightarrow \RR$ be a mapping.
    We call the \textbf{support} of $f$ the set
    $$\supp(f) \coloneq \overline{\{x \in X \mid f(x) \neq 0\}}.$$
    If $\supp(f)$ is a compact subset of $X$, then we say that $f$ is \textbf{compactly supported} (or of compact support).
    We denote by $C_c(X)$ the set of all continuous mappings $f : X \rightarrow \RR$ of compact support.
    It is a vector space over $\RR$.
\end{definitionenv}

\begin{propositionenv}\label{prop:5.3.4}
    Let $(X,\mathscr{T})$ be a locally compact Hausdorff space.
    \begin{enumerate}[label=(\arabic*)]
        \item Let $x \in X$, $F \subseteq X$ closed, $x \notin F$.
            There exists a compact neighborhood $C$ of $x$ such that $F \cap C = \varnothing$.
        \item Let $x \in X$, $V_x$ open neighborhood of $x$.
            There exists an open neighborhood $W_x$ of $x$ such that $\bar{W}_x \subseteq V_x$ and $\bar{W}_x$ is compact.
        \item Let $Y \subseteq X$ be compact, $V \subseteq X$ open such that $Y \subseteq V$.
            Then there exists an open $W$ such that $\bar{W}$ is compact
            $Y \subseteq W \subseteq \bar{W} \subseteq V.$
    \end{enumerate}
\end{propositionenv}
\begin{proofenv}
    \quad
    \begin{enumerate}[label=(\arabic*)]
        \item Let $N \ni x$ be a compact neighborhood of $x$.
            Then $N \cap F$ is compact.
            Since $X$ is Hausdorff,
            there exists open $U$ and $V$ such that $x \in U$, $N \cap F \subseteq V$, and $U \cap V = \varnothing$.
            Notice that $\bar{U} \subseteq N$, $\bar{U}$ is a closed subset of a compact set.
            So $\bar{U}$ is also compact.
            \begin{figure}[H]
                \begin{center}


                    % Pattern Info

                    \tikzset{
                    pattern size/.store in=\mcSize, 
                    pattern size = 5pt,
                    pattern thickness/.store in=\mcThickness, 
                    pattern thickness = 0.3pt,
                    pattern radius/.store in=\mcRadius, 
                    pattern radius = 1pt}
                    \makeatletter
                    \pgfutil@ifundefined{pgf@pattern@name@_4c78189ts}{
                    \pgfdeclarepatternformonly[\mcThickness,\mcSize]{_4c78189ts}
                    {\pgfqpoint{0pt}{0pt}}
                    {\pgfpoint{\mcSize+\mcThickness}{\mcSize+\mcThickness}}
                    {\pgfpoint{\mcSize}{\mcSize}}
                    {
                    \pgfsetcolor{\tikz@pattern@color}
                    \pgfsetlinewidth{\mcThickness}
                    \pgfpathmoveto{\pgfqpoint{0pt}{0pt}}
                    \pgfpathlineto{\pgfpoint{\mcSize+\mcThickness}{\mcSize+\mcThickness}}
                    \pgfusepath{stroke}
                    }}
                    \makeatother
                    \tikzset{every picture/.style={line width=0.75pt}} %set default line width to 0.75pt        

                    \begin{tikzpicture}[x=0.75pt,y=0.75pt,yscale=-1,xscale=1]
                    %uncomment if require: \path (0,300); %set diagram left start at 0, and has height of 300
                    
                    %Shape: Regular Polygon [id:dp5439473125317085] 
                    \draw  [line width=2.25]  (250.77,128.33) .. controls (239.5,128.33) and (228.23,128.12) .. (216.96,128.33) .. controls (209.27,128.47) and (195.48,137.96) .. (189.91,142.63) .. controls (187.41,144.73) and (183.86,145.55) .. (181.8,148.08) .. controls (179.21,151.24) and (177.2,154.83) .. (175.03,158.3) .. controls (171.2,164.42) and (167.44,188.16) .. (175.71,193.71) .. controls (194.61,206.41) and (213.93,201.93) .. (235.22,204.61) .. controls (254.7,207.07) and (272.22,212.2) .. (292.02,210.74) .. controls (308.44,209.53) and (322.51,193.23) .. (324.48,177.37) .. controls (325.22,171.45) and (326.9,166.09) .. (323.13,161.02) .. controls (315.39,150.62) and (301.57,151.37) .. (290.67,148.08) .. controls (282,145.46) and (279.35,142.81) .. (274.44,137.86) .. controls (272.23,135.64) and (260.2,127.76) .. (250.77,128.33) -- cycle ;
                    %Shape: Regular Polygon [id:dp7360409660603735] 
                    \draw  [line width=2.25]  (372.77,124.33) .. controls (361.5,124.33) and (350.23,124.12) .. (338.96,124.33) .. controls (331.27,124.47) and (317.48,133.96) .. (311.91,138.63) .. controls (309.41,140.73) and (305.86,141.55) .. (303.8,144.08) .. controls (301.21,147.24) and (299.2,150.83) .. (297.03,154.3) .. controls (293.2,160.42) and (289.44,184.16) .. (297.71,189.71) .. controls (316.61,202.41) and (335.93,197.93) .. (357.22,200.61) .. controls (376.7,203.07) and (394.22,208.2) .. (414.02,206.74) .. controls (430.44,205.53) and (444.51,189.23) .. (446.48,173.37) .. controls (447.22,167.45) and (448.9,162.09) .. (445.13,157.02) .. controls (437.39,146.62) and (423.57,147.37) .. (412.67,144.08) .. controls (404,141.46) and (401.35,138.81) .. (396.44,133.86) .. controls (394.23,131.64) and (382.2,123.76) .. (372.77,124.33) -- cycle ;
                    %Shape: Regular Polygon [id:dp28888133472491795] 
                    \draw  [line width=2.25]  (223.66,155.27) .. controls (218.18,155.27) and (212.7,155.19) .. (207.22,155.27) .. controls (203.48,155.33) and (196.78,159.32) .. (194.07,161.28) .. controls (192.86,162.16) and (191.13,162.51) .. (190.13,163.57) .. controls (188.87,164.9) and (187.89,166.41) .. (186.84,167.86) .. controls (184.97,170.43) and (183.14,180.41) .. (187.17,182.74) .. controls (196.36,188.07) and (205.75,186.19) .. (216.1,187.32) .. controls (225.57,188.35) and (234.09,190.51) .. (243.72,189.89) .. controls (251.7,189.38) and (258.55,182.54) .. (259.51,175.87) .. controls (259.86,173.39) and (260.68,171.13) .. (258.85,169.01) .. controls (255.08,164.64) and (248.37,164.95) .. (243.06,163.57) .. controls (238.85,162.47) and (237.56,161.36) .. (235.17,159.28) .. controls (234.1,158.34) and (228.25,155.03) .. (223.66,155.27) -- cycle ;
                    %Shape: Polygon Curved [id:ds5480576438816945] 
                    \draw  [line width=2.25] [line join = round][line cap = round] (317.23,138.71) .. controls (310.4,138.71) and (303.57,138.53) .. (296.75,138.71) .. controls (292.09,138.84) and (283.73,147.37) .. (280.36,151.56) .. controls (278.85,153.45) and (276.7,154.19) .. (275.45,156.46) .. controls (273.88,159.3) and (272.66,162.53) .. (271.35,165.64) .. controls (268.5,204) and (302.8,202.27) .. (307.81,207.26) .. controls (319.61,209.46) and (322.94,206.92) .. (334.94,205.61) .. controls (344.88,204.52) and (363.5,172) .. (341.4,156.46) .. controls (336.15,154.11) and (334.54,151.73) .. (331.57,147.28) .. controls (330.23,145.28) and (322.94,138.2) .. (317.23,138.71) -- cycle ;
                    %Shape: Free Drawing [id:dp5753203482155477] 
                    \draw  [line width=3] [line join = round][line cap = round] (239.09,176) .. controls (239.09,176.33) and (239.09,176.67) .. (239.09,177) ;
                    %Shape: Path Data [id:dp2275311698121395] 
                    \draw  [pattern=_4c78189ts,pattern size=6pt,pattern thickness=0.75pt,pattern radius=0pt, pattern color={rgb, 255:red, 0; green, 0; blue, 0}][line width=2.25]  (323.13,161.02) .. controls (326.9,166.09) and (325.22,171.45) .. (324.48,177.37) .. controls (323.61,184.38) and (320.37,191.49) .. (315.61,197.33) .. controls (309.57,195.97) and (303.62,193.68) .. (297.71,189.71) .. controls (289.44,184.16) and (293.2,160.42) .. (297.03,154.3) .. controls (297.91,152.91) and (298.75,151.49) .. (299.62,150.1) .. controls (308.35,151.74) and (317.4,153.33) .. (323.13,161.02) -- cycle ;
                    
                    % Text Node
                    \draw (192.13,166.97) node [anchor=north west][inner sep=0.75pt]    {$U$};
                    % Text Node
                    \draw (398,108.4) node [anchor=north west][inner sep=0.75pt]    {$F=\bar{F}$};
                    % Text Node
                    \draw (336.94,209.01) node [anchor=north west][inner sep=0.75pt]    {$V$};
                    % Text Node
                    \draw (225.07,164.68) node [anchor=north west][inner sep=0.75pt]    {$x$};
                    % Text Node
                    \draw (205.13,105.97) node [anchor=north west][inner sep=0.75pt]    {$N$};
                    
                    
                    \end{tikzpicture}
                \end{center}
            \end{figure}
            Moreover, $\bar{U} \cap F \subseteq N \cap F \subseteq V$, so $\bar{U} \cap F \subseteq \bar{U} \cap V = \varnothing$.
            We take $\bar{U}$ as the desired compact neighborhood of $x$.
        \item Let $F = X \setminus V_x$ closed.
            By (1), there exists a compact neighborhood $C$ of $x$ such that $C \cap F = \varnothing$.
            We take $W_x = C^\circ$ (the interior of $C$).
        \item For any $y \in Y$, let $W_y$ be an open neighborhood of $y$ such that $\bar{W}_y$ is compact.
            Since $Y$ is compact, there exists a finite set $\{y_1,\cdots,y_n\}$ such that
            $$Y \subseteq W \coloneq \bigcup_{i=1}^n W_{y_i}.$$
            Notice that
            $$\bar{W} = \bigcup_{i=1}^n \bar{W}_{y_i},$$
            hence $\bar{W}$ is compact and $Y \subseteq W \subseteq \bar{W} \subseteq V.$
    \end{enumerate}
\end{proofenv}

\begin{lemmaenv}[Urysohn's lemma]
    Let $(X,\mathscr{T})$ be a locally compact Hausdorff space,
    $K \subseteq X$ be compact, $U \subseteq X$ be open such that $K \subseteq U$.
    Then there exists a continuous mapping $f : X \rightarrow [0,1]$ continuous such that
    \begin{enumerate}
        \item For any $x \in K$, $f(x) = 1$.
        \item $f$ is compact support and $\supp(f) \subseteq U$.
    \end{enumerate}
\end{lemmaenv}
\begin{proofenv}
    By proposition \ref{prop:5.3.4}, there exists open $W$ such that $\bar{W}$ is compact and
    $$K\subseteq W \subseteq \bar{W} \subseteq U.$$
    We will construct a family $\left(U_r\right)_{r\in \QQ\cap[0,1]}$ of open subsets of $X$ such that
    \begin{enumerate}
        \item $K \subseteq U_1$.
        \item For any $r \in \QQ\cap[0,1]$, $\bar{U}_{r} \subseteq W$.
        \item For any $r,s \in \QQ \cap [0,1]$ with $r < s$, one has $U_r \subseteq \bar{U}_{s}$.
    \end{enumerate}
    For this purpose, we write $\QQ\cap[0,1]$ as a sequence $(r_n)_{n\in\NN}$ such that $r_0 = 1$.
    By proposition \ref{prop:5.3.4} (3), there exists an open $U_{r_0} \subseteq X$ such that
    $$K \subseteq U_{r_0} \subseteq \bar{U}_{r_0} \subseteq W.$$ 
    Suppose that $U_{r_0}, U_{r_1}, \ldots, U_{r_{n}}$ are constructed.
    We sort the elements $\{r_0,\cdots,r_n\}$ into a strictly decreasing sequence
    $$1 = s_0 > s_1 > \cdots > s_n.$$
    Here comes to a new element $r_{n+1}$.
    We insert $r_{n+1}$ into the sorted sequence, find out index $i \in \{0,\cdots,n\}$ such that $s_i > r_{n+1} > s_{i+1}$.
    We pick an open $U_{r_{n+1}}$ such that $\bar{U}_{r_{n+1}}$ is compact and
    $$\bar{U}_{s_i} \subseteq U_{r_{n+1}} \subseteq \bar{U}_{r_{n+1}} \subseteq U_{s_{i+1}}.$$
    If $s_n > r_{n+1}$, then let
    $$\bar{U}_{s_n} \subseteq U_{r_{n+1}} \subseteq \bar{U}_{r_{n+1}} \subseteq W.$$
    By induction, we construct the family $\left(U_r\right)_{r\in \QQ\cap[0,1]}$ satisfying the desired properties.

    \quad For $x \in X$, we let
    $$f(x) = \sup\{r \in \QQ\cap[0,1] \mid x \in U_r\}$$
    with the convention that $\sup \varnothing = 0$.
    Note
    $$\{x \in X \mid f(x) \neq 0\} = W.$$
    Hence $\supp(f) \subseteq W$ is compact.
    Also $\supp(f) \subseteq \bar{W} \subseteq U$ and since $K \subseteq U_1$, $f|_K = 1$.
    
    \quad It remains to check that $f$ is continuous.
    Let $t \in [0,1]$.
    By definition,
    $$\{x \in X \mid f(x) > t\} = \{x \in X \mid \exists r \in \QQ \cap \interval[open left] {t}{1}, x \in U_r\} = \bigcup_{r \in \QQ\cap \interval[open left] {t}{1}} U_r.$$
    Moreover, for any $t \in \interval[open left]{0}{1}$,
    $$\{x \in X \mid f(x) \ge t\} = \bigcap_{r \in \QQ\cap\interval[open]{0}{t}} \bar{U}_r.$$
    Indeed, if $f(x) \ge t$, then for any $r \in \QQ \cap \interval[open]{0}{t}$, there exists $s \in \QQ$ with $r < s < t$ such that $x \in U_s \subseteq \bar{U}_r$. 
    Conversely, if $x \in \bigcap_{r \in \QQ\cap\interval[open]{0}{t}} \bar{U}_r$, then for any $s < t$ we have $x \in \bar{U}_s$, which implies $f(x) \ge s$ for all $s < t$, hence $f(x) \ge t$.

    \quad Note that the set $\{x \in X \mid f(x) > t\}$ is equal to $\bigcup_{r \in \QQ\cap\interval[open]{t}{1}} U_r$, which is open. 
    Thus $f$ is continuous.


    \quad Let $I = \interval[open]{a}{b}$ be an open interval,
    where $a,b \in \RR$ with $a < b$.
    If $a \ge 0$ and $b > 1$,
    then 
    $$f^{-1}(I) = \{x \in X \mid f(x) > a\}.$$
    If $a < 0$ and $b \le 1$, then
    $$f^{-1}(I) = \{x \in X \mid f(x) < b\}.$$
    If $\{a,b\} \subseteq [0,1]$, then
    $$f^{-1}(I) = \{x \in X \mid f(x) > a\} \cap \{x \in X \mid f(x) < b\}.$$
    In all cases, $f^{-1}(I)$ is an open set of $X$.
    Since any open set in $X$ can be written as a union of intervals.
    We deduce that $f$ is continuous.
\end{proofenv}

\begin{corollaryenv}
    Let $(X,\mathscr{T})$ be a locally compact Hausdorff space.
    Then the topology $\mathscr{T}$ identifies with the coarsest topology that makes all the mappings in $C_c(X)$ continuous.
\end{corollaryenv}
\begin{proofenv}
    By definition, any mapping in $C_c(X)$ is continuous with respect to $\mathscr{T}$ on $X$.
    Let $\mathscr{T}'$ be a topology on $X$ such that any $f \in C_c(X)$ is continuous with respect to $\mathscr{T}'$.
    Let $U \neq \varnothing$ belong to $\mathscr{T}$.
    We want to prove that $U$ must belong to $\mathscr{T}'$.
    
    \quad For any $x \in X$, we pick a compact neighborhood $K_x$ of $x$ such that $K_x \subseteq U$.
    By Urysohn's lemma, there exists $f_x \in C_c(X)$ such that $\supp(f_x) \subseteq U$ and that
    $$K_x \subseteq V_x \coloneq \{y \in X \mid  f_x(y) > 0\}.$$
    So,
    $$U = \bigcup_{x \in U} K_x \subseteq \bigcup_{x \in U}V_x \subseteq U.$$
    
    \quad Thus, $U$ belongs to $\mathscr{T}'$.
\end{proofenv}

\begin{propositionenv}\label{prop:5.3.7}
    Let $(X,\mathscr{T})$ be a locally compact Hausdorff space which has a countable basis of $\mathscr{T}$.
    We claim there exist a sequence $\left(U_n\right)_{n \in \NN}$ of open subset of $X$ such that
    \begin{enumerate}
        \item For any $n \in \NN$, $\bar{U}_n$ is compact and contained in $U_{n+1}$.
        \item $\displaystyle \bigcup_{n \in \NN} U_n = X$
    \end{enumerate}
\end{propositionenv}
\begin{proofenv}
    Let $\mathscr{B}$ be a countable basis of $\mathscr{T}$ and $\mathscr{B}_c = \{V \in \mathscr{B} \mid \bar{V} \text{ compact}\}$ is also countable.
    We want to show that $\mathscr{B}_c$ also forms a basis of $\mathscr{T}$.
    For any $x \in U$, there exists an open $W_x \ni x$ such that $\bar{W}_x$ is compact and $W_x \subseteq \bar{W}_x \subseteq U$ (by \ref{prop:5.3.4}).
    Since $\mathscr{B}$ is a basis of $\mathscr{T}$, there exists $V_x \in \mathscr{B}$ such that $x \in V_x \subseteq W_x$.
    Note that $\bar{V}_x \subseteq \bar{W}_x$  and hence is compact.
    We obtain that $V_x \in \mathscr{B}_c$.
    Hence $\mathscr{B}_c$ forms a countable basis of $\mathscr{T}$.
    Write
    $$\mathscr{B}_c = \{V_n \mid n \in \NN\}.$$
    We construct $\left(U_n\right)_{n\in \NN}$ in a recursive way.
    First, $U_0 = V_0$.
    Suppose that $U_0,\cdots,U_n$ are such that $\bar{U}_{i}$ is compact, $V_i \subseteq U_i$, and
    $$\forall i \in \{0,\cdots,n-1\},\ \bar{U}_{i} \subseteq U_{i+1}.$$
    We want to construct $U_{n+1}$.
    Note that $\bar{U}_{n} \cup \bar{V}_{n+1}$ is compact,
    and by proposition \ref{prop:5.3.4} (3), there exists an open $U_{n+1}$ such that $\bar{U}_{n+1}$ is compact and 
    $$\bar{U}_{n} \cup \bar{V}_{n+1} \subseteq U_{n+1}.$$
    Note that
    $$X = \bigcup_{n \in \NN} V_{n} = \bigcup_{n \in \NN} U_{n} \subseteq X.$$
\end{proofenv}

\begin{lemmaenv}\label{lem:5.3.8}
    Let $X$ be a topological space with a countable topological basis.
    Any open cover $\{U_i\}_{i\in I}$ of $X$ has a countable subcover.
\end{lemmaenv}
\begin{proofenv}
    Let $(V_n)_{n\in \NN}$ be a family of open sets that forms a topological basis of $X$.
    Let $\{U_i\}_{i\in I}$ be any of open cover of $X$.
    For any $x \in X$, we choose $i(x) \in I$ such that $x \in U_{i(x)}$ and let $n(x) \in \NN$ be such that $x \in V_{n(x)}$.
    Let
    $$A = \{n(x) \mid x \in X\} \subseteq \NN,\quad X = \bigcup_{m \in A} V_{m}.$$
    For any $m \in A$, we pick some $x_m \in X$ such that $n(x_m) = m$.
    $$X = \bigcup_{m \in A} V_{m} \subseteq \bigcup_{m \in A} U_{i(x_m)}.$$
\end{proofenv}

\begin{propositionenv}\label{prop:5.3.9}
    Let $X$ be a non-empty locally compact Hausdorff space.
    Let $(F_n)_{n\in \NN}$ be a family of closed subsets of $X$.
    If 
    $$X = \bigcup_{n \in \NN} F_n,$$
    then there exists $N \in \NN$ such that $F_{N}^{\circ}$ is not empty.
\end{propositionenv}
\begin{proofenv}
    Suppose, by contradiction that $F^{\circ}_{n} = \varnothing$ for any $n \in \NN$.
    
    \quad Let $U_0$ be open, such that $\bar{U}_0$ is compact.
    Suppose that non-empty open subsets $U_0,\cdots,U_n$ have been chosen in a way that $\bar{U}_{i}$ is compact and contained in $U_{i-1}\setminus F_{i}$ (open).
    
    \quad We construct $U_{n+1}$ such that $\bar{U}_{n+1} \subseteq U_{n}\setminus F_{n+1}$.
    $(\bar{U}_{n})_{n \in \NN}$ is a decreasing sequence of compact subsets and hence $\displaystyle \bigcap_{n \in \NN} \bar{U}_{n} \neq \varnothing$.
    
    \quad This contradicts
    $$\varnothing \neq X = \bigcup_{n \in \NN} F_n.$$
    (Because $x \in \bigcap_{n \in \NN} U_n \Rightarrow x \notin F_{n}$.)
\end{proofenv}

\begin{propositionenv}
    Let $(G,*,\mathscr{T})$ be a topological group,
    $X$ be a topological space and $\varphi : G \times X \rightarrow X$ continuous left action which is \textit{transitive}.
    Assume that both $G$ and $X$ are locally compact with countable topological basis.
    Let $x_0 \in X$ and let $H = \stab(x_0)$.
    
    \quad Then the mapping
    $$\begin{array}{rrcl}
        f : & G/H & \rightarrow & X\\
        & a * H & \longmapsto & \varphi (a,x_0)
    \end{array}$$
    is a homeomorphism.
\end{propositionenv}
\begin{proofenv}
    By proposition \ref{prop:5.2.17}, $f$ is continuous bijection,
    so we only have to prove that $f$ is open.

    \quad Let $V \neq \varnothing$ be an open subset of $G$.
    Let $x \in f (V)$ and let $s \in G$ such that $f(s) = x$.
    By proposition \ref{prop:5.2.10} (1), there exists an open $U_{0} \ni e$,
    $$s * \iota(U_{0}) * U_{0} \subseteq V.$$
    Let $U \ni e$ be such that $\bar{U}$ is connect and $\bar{U} \subseteq U_0$.
    Since $G$ has a countable topological basis, then by lemma \ref{lem:5.3.8}, an open cover $(a * U)_{a \in G}$ has an countable subcover $(a_n * U)_{n \in \NN}$.
    Note that
    $$G = \bigcup_{n \in \NN} (a_n * U) = \bigcup_{n \in \NN} a_{n} * \bar{U}.$$
    We have that
    $$X = \bigcup_{n \in \NN} L_{a_{n}} \left(f(\bar{U})\right).$$
    Since $\bar{U}$ is compact, so is $L_{a_{n}} \left(f(\bar{U})\right)$.
    By proposition \ref{prop:5.3.9}, there exists $n \in \NN$ such that $L_{a_{n}}\left(f(U)\right)$ has a non-empty interior.
    We conclude that
    $$ W \coloneq f(\bar{U})^\circ \text{ is non-empty.}$$
    Let $y \in W$ and $b \in \bar{U}$ such that $f(b) = y$.
    Then
    $$\begin{aligned}
        & x = f(s) = L_{s * \iota(b)} = L_{s * \iota(b)}(y)\\
        \in & L_{s * \iota(b)}(W) \subseteq f(s * \iota(\bar{U}) * \bar{U}) \subseteq f(s * \iota(U_0) * U_0) \subseteq f(V).
    \end{aligned}$$
    This shows that $L_{s * \iota(b)}(W)$ is an open neighborhood of $x$ in $f(V)$.
    We conclude that $f(V)$ is open in $X$.
\end{proofenv}


\section{Radon Measure}


\begin{definitionenv}
    Let $(X,\mathscr{T})$ be a locally compact Hausdorff space such that has a countable topological basis.
    We call \textbf{positive linear functional} on $X$ any linear mapping $I : C_c (X) \rightarrow \RR$,
    such that for any $f \in C_c(X)$,
    if $\forall x \in X$, $f(x) \ge 0$, then $I(f) \ge 0$.
\end{definitionenv}

\begin{theoremenv}
    Let $(X,\mathscr{T})$ be a locally compact Hausdorff space such that has a countable topological basis and $I : C_c(X) \rightarrow \RR$ be a positive linear functional.
    
    \quad Then $I$ is an integral operator and determines a unique $\sigma$-finite Borel measure $\mu_{I}$ on $X$ such that
    $$\forall f \in C_c(X), \ \int_{X} f(x) \mu_{I} (\dd x) = I(f).$$
    The measure $\mu_{I}$ is called the \textbf{Radon measure} associated with $I$ on the Borel $\sigma$-algebra of $X$.
\end{theoremenv}
\begin{proofenv}
    Let $f, g \in C_{c}(X)$,
    then $\inf\{f,g\}$ is still continuous and
    $$\supp(\inf\{f,g\}) \subseteq \supp(f) \cup \supp(g).$$
    Hence, 
    $$\inf\{f,g\} \in C_c(X).$$
    i.e. $C_c(X)$ is a Riesz space.
    
    \quad Let $(f_n)_{n \in \NN}$ be a decreasing sequence in $C_c(X)$ that converges pointwise to the constant zero mapping.
    (We want to show that $\displaystyle \lim_{n \rightarrow \infty}I(f_n) = 0$ to show that $I$ is an integral operator.)
    Let $K$ be the support of $f_0$ ($\supp (f_n) \subseteq K$ for any $n$),
    $K$ is compact,
    so by the Urysohn's lemma,
    one can pick $f : X \rightarrow [0,1]$ such that $f \in C_c(X)$ and $f(x) = 1$ for any $x \in K$.
    Then for any $n \in \NN$,
    $$\forall x \in X, f_n(x) \le \|f_{n}\|_{L^\infty} f(x),$$
    where $\displaystyle \|f_{n}\|_{L^\infty} = \sup_{x \in X} |f_n(x)|$.
    Since $I$ is a positive linear functional,
    $$I(f_n) \le \|f_{n}\|_{L^\infty} I(f(x)).$$
    We want to check that
    $$\lim_{n \rightarrow \infty}I(f_n) = 0.$$
    Since the sequence $\left(\|f_{n}\|_{L^\infty}\right)_{n \in \NN}$ is decreasing and bounded from below (by $0$),
    it converges to some $c \in \RR$.

    \quad Assume by contradiction that $c > 0$.
    Then there exists $\varepsilon > 0$ such that $\|f_{n}\|_{L^\infty} > \varepsilon$ for all $n \in \NN$.
    Let
    $$K_{n} \coloneq \{x \in X \mid f_{n} (x) \ge \varepsilon\}$$
    The sequence $\left(K_{n}\right)_{n \in \NN}$ is a decreasing sequence of non-empty compact sets of $X$.
    Thus
    $$\bigcap_{n \in \NN} K_{n} \neq \varnothing.$$
    This contradicts, the assumption that $(f_n)_{n\in \NN}$ converges to the constant $0$ function pointwise.
    
    \quad To conclude, we want to use Carath\'eodory's (extension) theorem.
    (It gives us the existence of $\mu_I$.)

    \quad By proposition \ref{prop:5.3.7}, we have $(U_n)_{n \in \NN}$ family of open sets $\bar{U}_{i} \subseteq U_{i+1}$, $\bar{U}_{i}$ compact,
    and $\displaystyle X = \bigcup_{n \in \NN} U_{n}$.
\end{proofenv}

\begin{remark}
    Let $X$ be a locally compact Hausdorff space,
    $I : C_c(X) \rightarrow \RR$ be a positive linear functional.
    Then, for any $f : X \rightarrow \RR_{\ge 0}$ bounded Borel measurable mapping,
    one has $f \in \mathcal{L}^1(I)^\uparrow$.
    So $I(f)$ is well defined.
    \begin{enumerate}
        \item For any Borel set $A \subseteq X$, $\mathbbm{1}_{A} \in \mathcal{L}^1(I)^\uparrow$.
        \item Let $\mathcal{B}$  be the Borel $\sigma$-algebra, then 
            $$\begin{array}{rrcl}
                \mu_{I} : & \mathcal{B} & \rightarrow & \RR \cup \{+\infty\} \\
                & A & \longmapsto & I(\mathbbm{1}_{A})
            \end{array}$$
            is a measure on $(X,\mathcal{B})$, called a \textbf{Radon measure} associated to $I$.
    \end{enumerate}
\end{remark}

\begin{propositionenv}\label{prop:5.4.4}
    Let $X$ be a locally compact Hausdorff space having a countable topological basis,
    $I : C_{c} (X) \rightarrow \RR_{\ge 0}$ be a positive linear functional,
    $\mu_{I}$ be the corresponding Radon measure.
    \begin{enumerate}[label=(\arabic*)]
        \item For any $K \subseteq X$ compact, $\mu_I(K) < +\infty$.
        \item For any open subset $U$ of $X$, one has
        $$\mu_{I} (U) = \sup_{\substack{K \subseteq U\\K \text{ compact}}} \mu_{I}(K) = \sup_{\substack{(f : X \rightarrow [0,1]) \in C_c(X)\\ \supp(f) \subseteq U}} I(f).$$
        \item For any Borel subset $A$ of $X$,
        $$\mu_{I} (A) = \inf_{\substack{U \text{ open}\\ A \subseteq U}} \mu_{I} (U).$$
        \item For any Borel subset $A$ of $X$, one has 
        $$\mu_{I}(A) = \sup_{\substack{K \subseteq A\\K \text{ compact}}} \mu_{I}(K).$$
    \end{enumerate}
\end{propositionenv}
\begin{proofenv}
    \quad 
    \begin{enumerate}[label=(\arabic*)]
        \item There exists $f \in C_c (X)$, $f(X) \subseteq [0,1]$.
            such that $f|_K \equiv 1$, $\mathbbm{1}_{K} \le f$, so $\mu_{I} (K) \le I(f)$,
        \item There exists an increasing sequence $(K_n)_{n \in \NN}$ of compact subset of $U$ such that $\displaystyle \bigcup_{n \in \NN} K_{n} = U$.
            $$\mu_{I} (U) = \lim_{n \rightarrow +\infty} \mu_{I} (K_{n}) \le \sup_{\substack{K \subseteq U\\ K \text{ compact}}} \mu_{I} (K).$$
            For any compact $K \subseteq U$, there exists $f \in C_c(X)$ such that $f|_K \equiv 1$, $\supp(f) \subseteq U$.
            $\mathbbm{1}_{K} \le f \le \mathbbm{1}_{U}$, so $\mu_{I} (K) \le I(f) \le \mu_{I}(U)$.
            So,
            $$\sup_{\substack{K \subseteq U\\ K \text{ compact}}} \mu_{I} (K) \le \sup_{\substack{(f : X \rightarrow [0,1]) \in C_c(X)\\ \supp(f) \subseteq U}} I(f) \le \mu_{I} (U).$$
        \item We may assume that $\mathbbm{1}_{A} \in \mathcal{L}^1(I)$,
            $$\mu_{I}(A) = \inf_{\substack{f \in C_{c}(X)^\uparrow\\ f \ge  \mathbbm{1}_{A}}} I(f).$$
            ($\mu_{I}(A) < +\infty$.)

            Let $f \in C_{c}(X)^\uparrow$, $t \in \RR$.
            We claim that
            $$U_{f,t} \coloneq \{x \in X \mid f(x) > t\} = f^{-1} \left(\interval[open]{t}{+\infty}\right)$$
            is open.
            In fact, if $(f_n)_{n \in \NN}$ is an increasing sequence of function in $C_{c}(X)$ that converges pointwise to $f$,
            then
            $$U_{f,t} = \bigcup_{n \in \NN} \{x \in X \mid f_{n}(x) > t\}.$$
            Now let $(g_{n})_{n \in \NN}$ be a decreasing sequence in $C_c(X)^\uparrow$ such that $g_{n} \ge \mathbbm{1}_{A}$ for any $n \in \NN$ and $\dis \lim_{n \rightarrow +\infty} I(g_n) = \mu_{I}(A)$.
            Let $\varepsilon \in \interval[open]{0}{1}$.
            For any $n \in \NN$, 
            $$U_{g,1-\varepsilon} \coloneq \{x \in X \mid g_{n}(x) > 1 - \varepsilon\}$$
            is open, and $g_{n} \ge (1 - \varepsilon) \mathbbm{1}_{U_{g_{n}, 1 - \varepsilon}}$, $\mu_{I}(A) \le \mu_{I}(U_{g_n,1 - \varepsilon}) \le \frac{I(g_n)}{1 - \varepsilon}$.
            $\dis \inf_{\substack{U \supseteq A\\ U \text{ open}}} \mu_{I}(U) \le \frac{\mu_{I}(A)}{1 - \varepsilon}$.
            Take $\varepsilon \rightarrow 0$.
        \item Suppose that $\mu_{I}(A)$ is finite.
            Let $\varepsilon > 0$, and $U$ open such that $A \subseteq U$, and $\mu_{I}(U) \le \mu_{I}(A) + \varepsilon$.
            By (2), there exists a compact subset $F$ such that $F \subseteq U$, and $\mu_{I}(U\setminus F) \le \varepsilon$.
            By (3), there exists an open subset $V$ such that $U\setminus A \subseteq V$ and $\mu_{I}(V) \le 2 \varepsilon$.
            Let $K = F \setminus V$.
            It is compact and
            $$K \subseteq F \setminus (U \setminus A) = \{x \in F \mid x \notin (U \setminus A)\} \subseteq A.$$
            $$\mu_{I}(A \setminus K) \le \mu_{I}(U \setminus K) = \mu_{I}(U \setminus F) + \mu_{I}(F \setminus K) \le 3 \varepsilon.$$
            Suppose that $\mu_{I}(A) < +\infty$.
            Since $\mu_{I}$ is $\sigma$-finite,
            there exists an increasing sequence $(A_n)_{n \in \NN}$ of Borel subsets of $X$ such that $\dis A = \bigcup_{n \in \NN} A_{n}$ and $\mu(A_n) < +\infty$, for any $n \in \NN$.

            \quad For any $n \in \NN$, there exists a compact subset $K_n$ of $A_n$, such that
            $$\sup_{\substack{K\subseteq A\\K \text{ compact}}}\mu_{I}(K) \ge \mu_{I}(K_n) \ge \mu_{I}(A_n) - 1.$$
            Taking the limit when $n \rightarrow +\infty$.
            We get $\dis \sup_{\substack{K \subseteq A\\ K \text{ compact}}} \mu_{I}(K) = +\infty$.
    \end{enumerate}
\end{proofenv}


\section{Haar Measure}

We fix a locally compact Hausdorff topological group $(G,*,\mathscr{T})$ having a countable topological basis.
Denote by
$$C_{c}^{+} (G) \coloneq \{(f : G \rightarrow \RR_{\ge 0}) \in C_c(G) \mid f \not\equiv 0\}$$
Let
$$\begin{array}{rrcl}
    \| \cdot \|_{\sup} : & \cc(G) & \longrightarrow & \RR_{\ge 0}\\
    & f & \longmapsto & \dis \sup_{x \in G} |f(x)|.
\end{array}$$
For any $a \in G$, let $L_a : \cc(G) \rightarrow \cc(G)$ such that $L_{a} f (x) = f(a * x)$,
$R_a : \cc(G) \rightarrow \cc(G)$ such that $R_{a} f (x) = f(x * a)$.

Let
$$\begin{array}{rrcl}
    \iota : & \cc(G) & \longrightarrow & \cc(G)\\
    & f(\cdot) & \longmapsto & f (\iota (\cdot))
\end{array}$$

\begin{definitionenv}
    Let $I : \cc(G) \rightarrow \RR$ be a non-zero positive linear functional.
    If for any $a \in G$ and any $f \in \cc(G)$,
    $$I(L_a f) = I(f),$$
    we say that $I$ defines a \textbf{left Haar measure}.
    (Resp. $I(R_a f) = I(f)$ and right Haar measure.)
\end{definitionenv}

\begin{theoremenv}
    Let $(G,*,\mathscr{T})$ be a locally compact Hausdorff topological group.
    \begin{enumerate}[label=(\arabic*)]
        \item There exists a positive linear functional
            $$I : \cc(G) \rightarrow \RR$$
            that defines a left Haar measure.
        \item If $I_1$ and $I_2$ are two non-zero positive linear functionals $\cc(G) \rightarrow \RR$ that define left Haar measures,
            then there exists $\lambda > 0$ such that $I_2 = \lambda I_1$.
    \end{enumerate}
\end{theoremenv}


\section{Change of Variables}


\begin{definitionenv}
    Let $(G,*)$ be a group.
    If $(a,b) \in G \times G$, we define the \textbf{commutator} of $a$ and $b$ the element
    $$\commutator{a}{b}\coloneq a * b * \iota (a) * \iota(b).$$
    $$a * b = b * a \Leftrightarrow \commutator{a}{b} = e.$$
    Denote by $\commutator{G}{G}$ the subgroup of $G$ generated by all commutators.

    \quad If $g \in G$, denote by
    $$\begin{array}{rrcl}
        C_{g} : & G & \longrightarrow & G\\
        & x & \longmapsto & g * x * \iota(g)
    \end{array}$$
    the mapping of conjugate by $g$.
    (Check that $C_g$ is a group homomorphism.)
\end{definitionenv}

\begin{propositionenv}
    Let $a,b \in G$.
    \begin{enumerate}[label=(\arabic*)]
        \item $\iota(\commutator{a}{b}) = \commutator{b}{a}$.
        \item $C_g\left(\commutator{a}{b}\right) = \commutator{C_g(a)}{C_g(b)}$.
    \end{enumerate}
\end{propositionenv}
\begin{proofenv}
    \quad
    \begin{enumerate}[label=(\arabic*)]
        \item $$\begin{aligned}
            \iota(a * b * \iota(a) * \iota(b)) = & \iota(\iota (b)) * \iota(\iota(a)) * \iota(b) * \iota(a)\\
            = & b * a * \iota(b) * \iota(a)\\
        \end{aligned}$$
        \item $$C_g(a * b * \iota(a) * \iota(b)) = C_g(a) * C_g(b) * \iota(C_g(a)) * \iota(C_g(b))$$
            since $C_{g}$ is a group homomorphism.

    \end{enumerate}
\end{proofenv}

\begin{corollaryenv}
    $\commutator{G}{G}$ is a normal subgroup of $G$.
\end{corollaryenv}

\begin{definitionenv}
    Let $K$ be a field and $V$ be a finite-dimensional vector space over $K$
    $$\End_{K}(V) \coloneq \{K\text{-linear endomorphisms} V \rightarrow V\}.$$
    Let
    $$\GL(V) \coloneq \End_{K}(V)^\times = \{\varphi \in \End_{K}(V) \mid \det (\varphi) \neq 0\},$$
    called the \textbf{general linear group} of $V$.
    Let
    $$SL(V) \coloneq \{\varphi \in \End_{K}(V) \mid \det (\varphi) = 1\} = \ker (\det),$$
    called the \textbf{special linear group} of $V$.
\end{definitionenv}

\begin{lemmaenv}\label{lem:5.6.7}
    Any element of $SL(V)$ can be written as a composition of elements of the form $S_{i,j}(\lambda)$.
    Moreover, any $f \in \GL(V)$ can be decomposed as $f = g \circ h$,
    where $g \in SL(V)$, and
    $$\begin{array}{rrcl}
        h : & V & \longrightarrow & V\\
        & \dis \sum_{i=1}^n a_i e_i & \longmapsto & \dis a_1 \det (f) + \sum_{i=2}^n a_i e_i.
    \end{array}$$
\end{lemmaenv}
\begin{proofenv}
    Let $\Theta$ be the family of $f \in \GL(V)$ that can be written as $g \circ h$ with $g$ and $f$ satisfies the conditions in the statement.
    We will prove by induction on $n = \dim (V)$ that
    $$\Theta = \GL(V).$$
    The cases $n = 0$, $n = 1$ are trivial.
    Now we assume that $n \ge 2$, and the induction hypothesis holds for $W = \Span_{K}\{e_1,\cdots,e_{n-1}\}$
    and the basis $(e_i)_{i=1}^{n-1}$.
    Note that for any $F \in \GL(V)$ and $i\neq j$, $i,j \in \{1,\cdots,n\}$, $\lambda \in K$.
    $F \in \Theta$ is equivalent to $S_{i,j}(\lambda) \circ F \in \Theta$.
    So in the reasoning without loss of generality,
    we can replace $f$ by $S_{i,j}(\lambda) \circ f$.
    For any $f \in \GL(V)$ and $(i,j) \in \{1,\cdots,n\}$, let $c_{i,j}(f)$ be such that
    $$f(e_i) = \sum_{j=1}^{n} c_{i,j}(f) e_j.$$
    $f \in \GL(V)$ (it is invertible), so $c_{n,1}(f) = \cdots = c_{n,n}(f) = 0$ can't hold.
    
    \quad First consider the case when $c_{n,j}(f) \neq 0$, for some $j \in \{1,\cdots,n-1\}$.
    Put
    $$\lambda = \frac{1 - c_{n,n}(f)}{c_{n,j}(f)}.$$
    One has
    $$\left(S_{n,j}(\lambda)\circ f\right)(e_n) = f(e_n) + \lambda c_{n,j}(f) e_n = e_n + \sum_{i=1}^{n-1}c_{n,i}(f)e_{i}.$$
    In the case where $c_{1,1}(f) = \cdots = c_{n,n-1}(f) = 0$, by replacing $f$ by $S_{1,n}(1) \circ f$ to reduce the problem, to the previous case.
    Hence, without loss of generality, we may assume that $c_{n,n}(f) = 1$.
    Replacing $f$ by
    $$S_{1,n}(-c_{n,1}(f)) \circ \cdots \circ S_{n-1,n}(-c_{n,n-1}(f)) \circ f.$$
    We may assume that $c_{n,1}(f) \circ \cdots \circ c_{n,n-1}(f) = 0$.
    Let $I: W\rightarrow W$ be the automorphism defined as
    $$\bar{f}(\sum_{i=1}^{n-1}a_i e_i) = \sum_{i=1}^{n-1}a_i \sum_{j=1}^{n-1}c_{i,j}(f)e_j.$$
    Applying the induction hypothesis to $\bar{f}$, we reduce the problem to the case where
    $$f(e_n) = e_n,\quad f(e_1) = c_{1,1} (f) e_1 + c_{1,n} (f) e_n,$$
    $$\forall i \in \{2,\cdots,n-1\},\; f(e_i) = e_i + c_{i,n}(f) e_n.$$
    $$\begin{aligned}
        f(e_1) \wedge \cdots \wedge f(e_n) = & \left(c_{1,1}(f)e_1 + c_{1,n}(f) e_n\right) \wedge \left(e_2 + c_{2,n}(f) e_n\right) \wedge \cdots\\
            & \wedge \left(e_{n-1} + c_{n-1,n}(f) e_n\right) \wedge e_n\\
            & c_{1,1}(f) e_1 \wedge \cdots \wedge e_n.
    \end{aligned}$$
    So $\det(f) = c_{1,1} (f)$.
    Finally, the mapping
    $$S_{n,1}\left(-c_{1,n}(f)/c_{1,1}(f)\right) \circ S_{n,2}(-c_{2,n}(f)) \circ \cdots \circ S_{n,n-1}(-c_{n-1,n}(f)) \circ f$$
    sends $\dis \sum_{i=1}^{n}a_ie_i$ to $\dis a_1 \det(f) + \sum_{i=2}^{n}a_ie_i.$
\end{proofenv}

\begin{theoremenv}\label{thm:5.6.*}
    Let $K$ be a field and $V$ be a finite-dimensional vector space over $K$.
    Assume that, either $\dim_{K}(V) \neq 2$, or $K \neq \ZZ /(2\ZZ)$.
    Then
    $$[\GL(V), \GL(V)]_{\circ} = SL(V).$$
\end{theoremenv}
\begin{proofenv}
    Since
    $$\det : \GL(V) \rightarrow K^\times$$
    is a homeomorphism of groups from (non-commutative) $\GL(V)$ to (commutative) $K^\times$,
    so
    $$[\GL(V), \GL(V)]_{\circ} \subseteq \ker (\det) = SL(V).$$
    If $\dim (V) = 0$, then $\GL(V) = SL(V)$ are both groups with one element.
    Hence, $[\GL(V), \GL(V)]_{\circ} = SL(V)$.
    If $\dim (V) = 1$, then $\det : \GL(V) \rightarrow K^\times$ is an isomorphism, and
    $$SL(V) = [\GL(V), \GL(V)]_{\circ} = \{\id_{V}\}.$$
    In the rest of the proof, we fix a basis $(e_i)_{i=1}^{n}$ of $V$ over $K$ ($n = \dim (V)$).
    For $i \neq j$, $i,j \in \{1,\cdots,n\}$ and $\lambda \in K$, we denote by $S_{i,j}(\lambda)$ the $K$-linear mapping:
    $$\begin{array}{rrcl}
        S_{i,j}(\lambda) : & V & \longrightarrow & V\\
        & \dis x = \sum_{i=1}^n a_i e_i & \longmapsto & x + \lambda a_j e_i.
    \end{array}$$
    In other words,
    $$S_{i,j}(\lambda)(e_j) = e_{j} + \lambda e_{i}, \quad \forall k \in \{1,\cdots,n\}\setminus \{j\},\ S_{i,j}(\lambda)(e_{k}) = e_{k}.$$
    Note that $S_{i,j}(\lambda) \in SL(V)$, for any $i,j \in \{1,\cdots, n\}, \lambda \in K$.
    Indeed,
    $$\begin{aligned}
        S_{i,j}(\lambda)(e_1) \wedge \cdots \wedge S_{i,j}(\lambda)(e_n) = & e_1 \wedge \cdots \wedge e_{j-1} \wedge (e_{j} + \lambda e_i) \wedge e_{j + 1} \wedge \cdots \wedge e_{n}\\
        = & e_1 \wedge \cdots \wedge e_{n}.
    \end{aligned}$$
    By definition, $\left(S_{i,j}(\lambda)\right)^{-1} = S_{i,j} (-\lambda)$.
    For any bijection $\sigma : \{1,\cdots,n\} \rightarrow \{1,\cdots,n\}$, we put
    $$\begin{array}{rrcl}
        P_{\sigma} : & V & \longrightarrow & V\\
        & \dis \sum_{i=1}^n a_i e_i & \longmapsto & \dis \sum_{i=1}^n a_{\sigma(i)} e_i.
    \end{array}$$
    Notice that 
    $$\sum_{i=1}^{n} a_{\sigma(i)} e_i = \sum_{i=1}^{n} a_i e_{\sigma^{-1}(i)},$$
    One has
    $$P_{\sigma}^{-1} = P_{\sigma^{-1}}.$$
    Moreover, $P_{\sigma^{-1}} \circ S_{i,j}(\lambda) \circ P_{\sigma}$ sends $\dis \sum_{i=1}^{n}a_i e_i$ to $\dis \sum_{i=1}^{n}a_{\sigma(i)} e_{\sigma(i)} + \lambda a_{\sigma(j)} e_{\sigma(i)}$.
    Hence $P_{\sigma^{-1}} \circ S_{i,j}(\lambda) \circ P_{\sigma} = S_{\sigma(i),\sigma(j)} (\lambda).$
    For $b \in K^\times$, denote $D_b \in \GL(V)$ such that
    $$D_b(\sum_{i=1}^{n}a_ie_i) = b a_1 e_1 + \sum_{i=2}^{n}a_i e_i.$$
    If $n=2$ and $K\not \cong \ZZ/(2\ZZ)$, then there exists $b\in K\setminus\{0,1\}$.
    Then
    $$[S_{1,2}(1), D_{b}]_{\circ} = S_{1,2}(1) \circ D_{b} \circ S_{1,2}(-1) \circ D_{b - 1}.$$
    $$\begin{aligned}
        [S_{1,2}(1), D_{b}]_{\circ} (a_1e_1 + a_2e_2) = & \left(S_{1,2}(1) \circ D_{b} \circ S_{1,2}(-1)\right)\left(a_1 b^{-1} e_1 + a_2 b_2\right)\\
        = & S_{1,2}(1) \circ D_{b} \left(a_1 b^{-1} e_1 + a_2 e_2 - a_2 e_1\right)\\
        = & S_{1,2}(-1) ((a_1 - b a_2) e_1 + a_2 e_2)\\
        = & (a_1 + (1 - b)a_2)e_1 + a_2 e_2\\
        = & S_{1,2}(1 - b) (a_1 e_1 + a_2 e_2).
    \end{aligned}$$
    Hence $$[S_{1,2}(1), D_{b}]_{\circ} = S_{1,2}(1 - b).$$
    If $n \ge 2$, then
    $$\begin{aligned}
        [S_{1,3}(1), S_{3,2}(1)]_{\circ} (x) = & S_{1,3}(1) \circ S_{3,2}(1) \circ S_{1,3}(-1) \circ S_{3,2}(-1)(x)\\
        = & S_{1,3}(1) \circ S_{3,2}(1) \circ S_{1,3}(-1) (x - a_2 e_3)\\
        = & S_{1,3}(1) \circ S_{3,2}(1) (x - a_2 e_3 - (a_3 -a_2)e_1)\\
        = & S_{1,3}(1 - a_2 - a_3 - (a_3 -a_2)e_2 +a_2 e_3)\\
        = & x - (a_3 - a_2) e_1 + a_3 e_1\\
        = & x + a_2 e_1\\
        = & S_{1,2}(1)(x).
    \end{aligned}$$
    In both cases $n = 2$ and $n > 2$, we obtained $t \in K\times$ such that
    $$S_{1,2}(t) \in [\GL(V), \GL(V)]_{\circ}.$$
    For general $\lambda \in K^\times$, one has
    $$S_{1,2}(\lambda) = D_{\lambda t^{-1}} \circ S_{12}(t) \circ D_{\lambda^{-1} t}.$$
    Since
    $$\begin{aligned}
        D_{\lambda t^{-1}} \circ S_{1,2}(t) \circ D_{\lambda^{-1} t} (\sum_{i=1}^{n}a_ie_i) = & D_{\lambda t^{-1}} \circ S_{1,2}(t) = a_1 \lambda^{-1} t + \sum_{i=2}^{n}a_i e_i\\
        = & D_{\lambda t^{-1}} \left(a_1 \lambda^{-1} t e_1 + a_2 t e_1 + \sum_{i=2}^{n}a_i e_i\right)\\
        = & S_{1,2}(\lambda) \left(\sum_{i=1}^{n}a_i e_i\right).
    \end{aligned}$$
    Finally for a general $(i,j) \in \{1,\cdots,n\}^2$, $i \neq j$, there exists a permutation $\sigma : \{1,\cdots,n\} \rightarrow \{1,\cdots,n\}$ such that $\sigma(1) = i$, $\sigma(2) = j$.
    We have $P_{\sigma^{-1}} \circ S_{i,j}(\lambda) \circ P_{\sigma} = P_{1,2}(\lambda)$.
    By lemma \ref{lem:5.6.7},
    $$[\GL(V), \GL(V)]_{\circ} = SL(V).$$
\end{proofenv}

\begin{remark}
    $$\det : \GL(V) \rightarrow K^\times.$$
    For any $\tau \in [\GL(V), \GL(V)]_{\circ}$, $\det(\tau) = 1$.
    So $\tau \in SL(V)$.
\end{remark}

\begin{remark}[Some assumptions and settings in \ref{thm:5.6.*}]\label{5.6.8}
    Let $\varphi : \GL(V) \rightarrow K^\times$ be a homeomorphism of groups.
    There exists an endomorphism of groups $\psi : K^\times \rightarrow K^\times$ such that
    $$\varphi = \psi \circ \det.$$
    \quad In fact, since $SL(V) = [\GL(V), \GL(V)]_{\circ}$, $\varphi$ sends elements of $SL(V)$ to $1$.
    Fix basis $(e_i)_{i=1}^{n}$ of $V$ over $K$.
    Let $D_{b} : V \rightarrow V$,
    $$\begin{array}{rrcl}
        \psi : & K^\times & \longrightarrow & K^\times\\
        & b & \longmapsto & \varphi(D_{b}).
    \end{array}$$
    This is also a homomorphism of groups.
    $$\psi(bb') = \psi(D_{bb'}) = \varphi(D_{b}\circ D_{b'}) = \varphi(D_{b}) \varphi(D_{b'}) = \psi(b)\psi(b').$$
    By lemma \ref{lem:5.6.7}, for any $f \in \GL(V)$, there exists $g \in SL(V)$ such that
    $$f = g \circ D_{\det(f)}.$$
    Hence $\varphi(f) = \varphi(g) \varphi(D_{\det(f)}) = \psi(\det(f)).$
\end{remark}

\begin{theoremenv}[Linear change of variables]\label{thm:5.6.9}
    Let $V$ be a finite-dimensional vector space over $\RR$ and $\nu$ be a Haar measure on $V$.
    For any $f \in \GL(V)$, one has
    $$f_*(\nu) = \left| \det(f)\right|^{-1} \nu.$$
    In other words, for any Borel subset $A \subseteq V$,
    $$\nu(f(A)) = \left| \det(f)\right| \nu(A).$$
\end{theoremenv}
\begin{proofenv}
    Let $(e_i)_{i=1}^n$ be a basis of $V$ over $\RR$.
    $$\begin{array}{rrcl}
        g : & \RR^n  & \longrightarrow & V\\
        & (a_1,\cdots,a_n) & \longmapsto & \sum_{i=1}^{n} a_i e_i.
    \end{array}$$
    Note that, $g$ is an isomorphism of topological groups (isomorphism of groups is also a homeomorphism).
    Let $\eta$ be the Lebesgue measure on $\RR^n$.
    Then $\eta^{\otimes n}$ is a Haar measure on $\RR^n$,
    and $g_*(\eta^{\otimes n})$ is a Haar measure on $V$.
    Without loss of generality, $\nu = g_*(\eta^{\otimes n})$.
    For $u \in \cc(V)$ and $y \in V$, one has
    $$\begin{aligned}
        \int_{V} u(x + y)f_*(\nu)(\dd x) = & \int_{V} u(f(x) + y) \nu(\dd x) = \int_{V} u(f(x + f^{-1}(y))) \nu(\dd x)\\
        = & \int_{V} u(f(x)) \nu(\dd x) = \int_{V} u(x)f_*(\nu)(\dd x).
    \end{aligned}$$
    We conclude $f_*(\nu)$ is a Haar measure on $V$.
    There exists $\phi(f) \in \RR_{>0}$ such that $f_*(\nu) = \varphi(f) \cdot \nu$.
    Moreover, for any $f_1,f_2 \in \GL(V)$,
    $$\left(f_2 \circ f_1\right)_* (\nu) = \left(f_2\right)_*\left((f_1)_*(\nu)\right) = \varphi(f_2) \cdot \varphi(f_1) \cdot \nu.$$
    Hence $\varphi$ is a morphism of groups from $\GL(V)$ to $\RR_{>0}$
    By \ref{5.6.8}, there exists a group homomorphism $\psi : \RR^\times \rightarrow \RR_{>0}$ such that
    $$\varphi = \psi \circ \det.$$
    In particular, if $f \in SL(V)$, then $\varphi(f) = 1$.
    Let $b \in V$.
    For any $u \in \cc(V)$, one has 
    $$\begin{aligned}
        \int_{V} u(x) (D_{b})_* (g_*(\eta^{\otimes n})) = & \int_{\RR^n} u(D_b(g(t_1,\cdots,t_n)))\eta^{\otimes n} \left(\dd (t_1,\cdots,t_n)\right)\\
        = & \int_{\RR^n} u(bt_1e_1 + t_2e_2 + \cdots + t_n e_n)\eta^{\otimes n}(\dd (t_1,\cdots,t_n))\\
        = & \int_{\RR^{n-1}}\left(\int_{\RR} u(bt_1e_1 + \sum_{i=2}^{n}t_ie_i)\eta(\dd t_1)\right) \eta^{\otimes n-1} \left(\dd (t_2,\cdots,t_n)\right).
    \end{aligned}$$
\end{proofenv}

\begin{lemmaenv}
    Let $h : \RR \rightarrow \RR$ be a non-negative Borel mapping.
    For any $b \in \RR^\times$, one has
    $$\int_{\RR} h(bt) \dd t = |b|^{-1} \int_{\RR} h(t) \dd t.$$
\end{lemmaenv}
\begin{proofenv}
    The above arguments applied to the additive group $\RR$ shows that there exists $\lambda : \RR^\times \rightarrow \RR_{>0}$ such that
    $$\int_{\RR} h(bt) \dd t = \lambda(b) \int_{\RR} h(t) \dd t.$$
    For any $b \in \RR^\times$, and non-negative Borel mapping $h$.
    Fixing $h \in \ccp(\RR)$, by dominated convergence,
    we obtain that $\lambda : \RR^\times \rightarrow \RR_{>0}$ is a continuous.
    By corollary \ref{cor:5.2.26}, there exists $\theta \in \RR$ such that $\lambda(b) = |b|^{\theta}$.
    Let's consider $b = 2$ and $h = \mathbbm{1}_{\interval[open right]{0}{2}}$,
    $$\int_{\RR} \mathbbm{1}_{\interval[open right]{0}{1}}(t) \dd t = \lambda(2) \int_{\RR} \mathbbm{1}_{\interval[open right]{0}{2}}(t) \dd t.$$
    So that $\lambda(2) = \frac{1}{2}$, so
    $$\forall b \in \RR,\ \lambda(b) = |b|^{-1}.$$
    By lemma, we get
    $$\int_{V} u(x) (D_b)_*(g_*(\eta^{\otimes n})) = |b|^{-1} \int_{V} u(x) g_{*} (\eta^{\otimes n})(\dd x).$$
    Therefore, by lemma \ref{lem:5.6.7}, we obtain
    $$\forall f \in \GL(V),\ \psi(f) = |\det(f)|^{-1}$$
    as required.
\end{proofenv}

\begin{theoremenv}[Change of variables]
    Let $E$ be a finite-dimensional vector space over $\RR$,
    $\nu$ be a Haar measure on $E$.
    Let $U,V$ be open subsets of $E$ and $\varphi : U \rightarrow V$ be a $C^1$-diffeomorphism.
    Then, for any Borel measurable mapping $f : V \rightarrow \RR_{\ge 0}$,
    one has
    $$\int_{V} f(y) \nu(\dd y) = \int_{U} f(\varphi(x)) |\det(\DD \varphi(x))| \nu(\dd x).$$
\end{theoremenv}
\begin{proofenv}
    Without loss of generality, we assume that $E = \RR^n$, $n = \dim(E)$.
    The case $n = 0$ is trivial, so assume that $n \ge 1$.
    We equip $\RR^n$ with the sup norm $\|\cdot\|$.
    Let $\Theta$ be the set hyperrectangles, i.e. the sets of the form:
    $$A = \prod_{i=1}^n \interval[open left]{a_i}{b_i}$$
    such that $\{a_1,b_1,a_n,b_n\} \subseteq \QQ$ and $\bar{A} \subseteq U$.
    If $b_1 - a_1 = \cdots = b_n - a_n$, then we say that $A$ is a hypercube.
    In general, $b_i - a_i$, $i \in \{1,\cdots,n\}$ are called $i$-th side length of the hyperrectangle.
\end{proofenv}

\begin{lemmaenv}\label{lem:5.6.12}
    For any hypercube $A \in \Theta$, one has
    $$\nu\left(\varphi(A)\right) \le \sup_{x \in A} \|D(\varphi(x))\|^n \nu(A).$$
\end{lemmaenv}
\begin{proofenv}
    Let $l$ be the side length of the hypercube.
    By the mean value inequality,
    for any $y,y' \in A$,
    $$\|\varphi(y) - \varphi(y')\| \le (\sup_{x \in A}\|\DD \varphi(x)\|)\|y-y'\|.$$
    So $\varphi(A)$ is contained in the hypercube of side length
    $$\left(\sup_{x \in A}\|D\varphi(A)\|\right)l.$$
\end{proofenv}

\begin{lemmaenv}
    For any hyperrectangle $A \in \Theta$, one has
    $$\nu(\varphi(A)) \le \int_{A} |\det(\DD \varphi(x))| \nu(\dd x).$$
\end{lemmaenv}
\begin{proofenv}
    Let $\varepsilon > 0$.
    Since the vertices of $A$ have rational coordinates,
    we can decompose $A$ into \textit{disjoint} finite union of hypercubes $(A_i)_{i\in I}$, $I$ finite, $\forall i \in I$, $A_i \in \Theta$,
    such that, for any $i \in I$ and any $x,y \in A_i$, one has
    $$\|\DD \varphi(y)^{-1} \circ \DD \varphi(x)\| \le (1 + \varepsilon)^{\frac{1}{n}}.$$
    This comes from the fact that $\DD \varphi$ is continuous,
    hence uniformly continuous on the compact set $\bar{A} \subseteq U$.
    Then
    $$\begin{aligned}
        \|\DD \varphi(y)^{-1} \circ \DD \varphi(x)\| \le & 1 + \|\DD \varphi(y)^{-1} \circ \DD \varphi(x) - \id_{\RR^n}\|\\
        \le & 1 + \|\DD \varphi(y)^{-1} \| \cdot \| \DD \varphi(x) - \DD \varphi(y)\|.
    \end{aligned}$$
    By the intermediate value theorem,
    for any $i \in I$, there exists $x_i \in A_i$ such that 
    $$|\det(\DD \varphi(x_i))| \nu(A_i) = \int_{A_i} |\det(\DD \varphi(x))| \nu(\dd x).$$
    By lemma \ref{lem:5.6.12}, 
    $$\nu(\DD(\varphi(x_i)^{-1})(\varphi(A_i))) \le \sup_{x \in A_i}\|\DD \varphi(x)^{-1} \circ \DD\varphi(x)\|^n \nu(A_i) \le (1 + \varepsilon)\nu(A_i).$$
    Since
    $$\varphi(A_i) = \DD\varphi (x_i) \left(\DD \varphi(x_i)^{-1}(\varphi(A_i))\right).$$
    By theorem \ref{thm:5.6.9},
    $$\begin{aligned}
        \nu(\varphi(A_i)) = & |\det(\DD \varphi(x_i))| |\nu (\DD(\varphi(x_i)^{-1})(\varphi(A_i)))|\\
        \le & |\det(\DD \varphi(x_i))| \cdot (1 + \varepsilon) \cdot \nu(A_i)\\
        = & (1 + \varepsilon) \int_{A_i} |\det(\DD \varphi(x))| \nu(\dd x).
    \end{aligned}$$
    Taking the union with respect to $i \in I$,
    we obtain
    $$\nu(\varphi(A)) \le (1 + \varepsilon) \int_{A} |\det(\DD \varphi(x))| \nu(\dd x).$$
    Note that $\Theta$ is countable and forms a (set theoretic) semi-ring on $U$.
    For any non-empty open set $U' \subseteq U$, one has
    $$U' = \bigcup_{\substack{A \in \Theta\\A \subseteq U'}}A.$$
    In particular, $U'$ can be written as a disjoint union of a countable family of elements of $\Theta$.
    By the monotone convergence theorem of Beppo-Levi, we obtain that,
    for any open, non-empty $U'\subseteq U$,
    $$\nu(\varphi(U')) \le \int_{U'} |\det(\DD \varphi(x))| \nu(\dd x).$$
    We notice that both $\nu \circ \varphi$ and $A \mapsto \int_{A} |\det(\DD \varphi(x))| \nu(\dd x)$ are Radon on $U$.
    By \ref{prop:5.4.4}, we obtain that for any Borel set $A$,
    $$\nu(\varphi(A)) \le \int_{A} |\det(\DD \varphi(x))|\nu(\dd x).$$
    Therefore, for any non-negative Borel mapping $f : V \rightarrow \RR$, one has
    $$\int_{V} f(y) \nu(\dd y) \le \int_{U} f(\varphi(x)) |\det(\DD \varphi(x))|\nu(\dd x)$$
    Indeed, we can write $f$ as a pointwise limit of the increasing sequence
    $$f_n = \sum_{k = 0}^{n2^n - 1}\frac{k}{2^n}\mathbbm{1}_{f^{-1}(\interval[open right]{\frac{k}{2^n}}{\frac{k + 1}{2^n}})}.$$
    Applying $\nu(\varphi(A)) \le \int_{A} |\det(\DD \varphi(x))|\nu(\dd x)$, we obtain
    $$\begin{aligned}
        \int_{V} f_n(y) \nu(\dd y) = & \sum_{k = 0}^{n2^n} \frac{k}{2^n} \nu(\varphi(f^{-1}(\interval[open right]{\frac{k}{2^n}}{\frac{k + 1}{2^n}})))\\
        \le & \nu(\varphi(A)) \le \int_{A} |\det(\DD \varphi(x))|\nu(\dd x)\\
        \le & \sum_{k=0}^{n2^n} \frac{k}{2^n} \int_{\varphi^{-1}(f^{-1}\interval[open right]{\frac{k}{2^n}}{\frac{k + 1}{2^n}})} |\det(\DD \varphi(x))|\nu(\dd x)\\
        = & \int_{U} f_{n} (\varphi(x)) |\det(\DD \varphi(x))|\nu(\dd x).
    \end{aligned}$$
    Taking the monotone limit $n \to +\infty$,
    we get
    $$\int_{V} f(y) \nu(\dd y) \le \int_{U} f(\varphi(x)) |\det(\DD \varphi(x))|\nu(\dd x).$$
    Applying this to $\varphi^{-1}$ and $h(x) = f(\varphi(x))|\det(\varphi(x))|$, we obtain
    $$\int_{U} f(\varphi(x)) |\det(\DD \varphi(x))|\nu(\dd x) \le \int_{V} f(y) |\det(\DD \varphi(\varphi^{-1}(y)))||\det \DD \varphi^{-1}(y)|\nu(\dd y).$$
    Since 
    $$\DD \varphi^{-1} (y) = \DD \varphi(\varphi^{-1}(y))^{-1},$$
    we obtain,
    $$\int_{U} f(\varphi(x)) |\det(\DD \varphi(x))|\nu(\dd x) \le \int_{V} f(y) \nu(\dd y).$$
    Hence the equality holds.
\end{proofenv}
